 \documentclass[12pt,twoside]{amsart} 
\usepackage[papersize={210mm,297mm},left=15mm,right=15mm,top=20mm,bottom=20mm]{geometry}
\usepackage{graphicx}
\usepackage{mathrsfs}
\usepackage{color}
\usepackage{amssymb}
\usepackage{amsthm}
\usepackage{xcolor}
\definecolor{slategray}{RGB}{112,138,144}


\usepackage[colorlinks=true]{hyperref}
\hypersetup{linkcolor=slategray,urlcolor=slategray,citecolor=slategray}
\usepackage[preserveurlmacro]{breakurl}

\newcommand{\tst}{\textstyle}

\newcommand{\E}{{\mathbb{E}}}
\newcommand{\I}{{\rm i}}
\newcommand{\dd}{{\rm d}}

%\newcommand{\pni}{\par\noindent}
%\newcommand{\pmni}{\par\medskip\noindent}
%\newcommand{\pbni}{\par\bigskip\noindent}

\newcommand{\beq}{\begin{equation}}
\newcommand{\eeq}{\end{equation}}
\newcommand{\neeq}{\nonumber\end{equation}}

\newcommand{\sfrac}[2]{{\textstyle \frac{#1}{#2}}}
\newcommand{\dbar}{{\mkern2mu\mathchar'26\mkern-2mu\mkern-9mud}}
\newcommand{\tr}{{\rm tr }}
\newcommand{\del}{\partial}

\parindent=0pt
\parskip=5pt

\renewcommand\Re{\operatorname{Re}}
\renewcommand\Im{\operatorname{Im}}

%\newcommand{\cB}{{\mathcal B}}
%\newcommand{\cE}{{\mathcal E}}
%\newcommand{\cG}{{\mathcal G}}
%\newcommand{\cH}{{\mathcal H}}
%\newcommand{\cS}{{\mathcal S}}
%\newcommand{\cT}{{\mathcal T}}
%
%\newcommand{\scD}{{\mathscr D}}
%\newcommand{\scE}{{\mathscr E}}
%\newcommand{\scF}{{\mathscr F}}
%\newcommand{\scX}{{\mathscr X}}
%
%\newcommand{\bC}{{\mathbb C}}
%\newcommand{\bM}{{\mathbb M}}
%\newcommand{\bN}{{\mathbb N}}
%\newcommand{\bR}{{\mathbb R}}
%\newcommand{\bZ}{{\mathbb Z}}

\newcommand{\al}{\alpha}
\newcommand{\de}{\delta}
\newcommand{\et}{\eta}
\newcommand{\ga}{\gamma}
\newcommand{\si}{\sigma}
\newcommand{\ph}{\phi}
\newcommand{\ch}{\chi}
\newcommand{\ps}{\psi}
\newcommand{\om}{\omega}

\newcommand{\etq}{{\bar\eta}}
\newcommand{\phq}{{\bar\phi}}
\newcommand{\chq}{{\bar\chi}}
\newcommand{\psq}{{\bar\psi}}

\newcommand{\veps}{\varepsilon}

\newcommand{\cP}{{\mathcal{P}}}
%\newcommand{\Ga}{\Gamma}
%\newcommand{\Ps}{\Psi}
%\newcommand{\Om}{\Omega}
%\newcommand{\Tr}{{\rm tr}}


\usepackage{graphicx}
\usepackage{appendix}
\usepackage{cases}

\newtheorem{theo}{\textsc{Theorem}}[section]
\newtheorem{lem}{\textsc{Lemma}}[section]
\newtheorem{defi}{\textsc{Definition}}[section]
\newtheorem{coro}{\textsc{Corollary}}[section]
\newtheorem{dem}{\textsc{Proof}}
\newtheorem{prop}{\textsc{Proposition}} [section]
\newtheorem{rem}{\textsc{Remark}}[section]
\newtheorem{exe}{\textsc{Example}}[section]
\numberwithin{equation}{section}
\DeclareMathOperator{\pp}{{\bf{p}}}
\DeclareMathOperator{\xx}{{\bf{x}}}
\DeclareMathOperator{\yy}{{\bf{y}}}
\DeclareMathOperator{\zz}{{\bf{z}}}
\DeclareMathOperator{\qq}{{\bf{q}}}
\DeclareMathOperator{\kk}{{\bf{k}}}
\newcommand*{\FD}{\hfill\ensuremath{\blacksquare}}

\newcommand{\reftitel}[1]{{\sl #1}}
\newcommand{\enquote}[1]{``#1''}
\newcommand{\neu}[1]{{\color{blue}#1}}
\newcommand{\alt}[1]{{\color{brown}#1}}


\begin{document}

\title[]{The majorant method for the fermionic effective action}
\author{Wilhelm Kroschinsky, Domingos H. U. Marchetti, Manfred Salmhofer} 
\address{
Institute for Applied Mathematics, Hausdorff-Center for Mathematics, University of Bonn, Endenicher Allee 60, 53115 Bonn, Germany
 --- %\\[1ex]
Instituto de F\' isica, Universidade de São Paulo, Brazil
 --- %\\[1ex]
Institut f\" ur Theoretische Physik, Universit\" at Heidelberg,
Philosophenweg~19, 69120 Heidelberg, Germany}

\email{kroschinsky@iam.uni-bonn.de,marchett@if.usp.br,salmhofer@uni-heidelberg.de}

\date{\today}

\begin{abstract}
\noindent
We revisit the problem of controlling Polchinski's equation by the solution of an associate Hamilton-Jacobi equation which determines a norm majorant for the fermionic effective action. This method, referred to as the Majorant Method, was first introduced by D. Brydges and J. Wright in 1988, but its original formulation contains a gap which has never been addressed. We overcome this gap and show that the majorant equation and its existence condition are analogous to the ones originally obtained by Brydges and Wright. As an application of the method, we investigate systems with quartic perturbations. % and special attention is given to the $\phi^{4}_{4}$ theory. 
\end{abstract}

\maketitle

\section{Introduction}

Polchinski's equation \cite{Polch} has revolutionized proofs of perturbative renormalizability, by shifting the emphasis in proofs from individual graphs to the Green functions of the quantum field theory. The structure of Polchinski's equation is that of a nonlinear heat equation in field space, which is infinite-dimensional in a continuum theory and very high-dimensional in most regularizations, such as the lattice regularization. 

Starting with \cite{KKS}, his method has been developed to the point where almost all results in perturbative renormalization have been reproven \cite{KeKoQED,KeKoZi,KoMue}, and many have been extended significantly. \cite{HoHo,KoHoOPE,HoStM,Sa98,MSbook}. 
The simplicity of the strategy of proof makes one want to use it also nonperturbatively, i.e.\ at small coupling, but not in the framework of formal power series. One case where this has been done is in the field-theoretical formulation of the KAM problem \cite{DeSiKupi}, which is, however, the case where the loops of quantum field theory are absent. 

Pioneering work relating Polchinski's equation and well-established methods of constructive quantum field theory was first done by Brydges and Kennedy. They showed that tree and polymer expansions can be obtained as explicit (convergent) series solutions of the differential equation, and that one can give a majorant for the solution of Polchinski's equation that satisfies a Hamilton-Jacobi equation. In a further paper, Brydges and Wright studied this majorant for fermionic systems \cite{BryWri}. It was later discovered that the proof in \cite{BryWri} contains a gap \cite{BryWriErr}, and it has since remained an open problem to complete the proof. This is the first purpose of the present paper. 

In the mean time, much work was devoted to simplifying and extending the fermionic tree expansions. While many works rely on the Brydges-Battle-Federbush interpolation, one can also use Polchinski's equation to generate an interpolation formula that has the necessary positivity properties. Continuous scale decompositions and flows were also used to construct the Gross-Neveu model in \cite{DR-GN} and two-dimensional many-fermion systems in \cite{DR1,DR2}. 

These constructions are, however, not as simple as the original Polchinski proof -- partly this is an inherent difference between perturbative and non-perturbative methods, but the strategy appears closer to that of tree expansions. The second purpose of this paper is to indicate how Polchinski-type estimates may be done using the majorant method. This is only an exemplification corresponding to the restriction to `completely convergent graphs' in other approaches; a full construction of fermionic models is possible along these lines, but is deferred to future work. 

Let us say a few words about how this paper is organized. For the reader's convenience, we start by recalling some standard definitions and results from the theory of Grassmann algebra. This is done in Section \ref{grassmannoverview}. In Section \ref{norms}, we quickly introduce the relevant norms to be used on the Grassmann algebra, which are in accordance with Brydges and Wright's paper. Section \ref{rgtransformation} is devoted to introducing a the renormalization group transformation and its connection to Polchinsky's equation. Our main result, Theorem \ref{majorant}, is stated and proved in Section \ref{majorantmethod}, where the Majorant Method is developed. Finally, some applications of the method using quartic interactions as initial condition are then discussed in Section \ref{quarticapplications}. 

Note added: After completion of this paper we learned about the recent work of Pawel Duch \cite{Duch}, who constructs the two-dimensional Gross-Neveu model using the Polchinski equation. 


\section{Grassmann Algebras - An Overview}\label{grassmannoverview}

Fermionic systems are naturally described in terms of Grassmann algebras \cite{berezin2012method}. For this reason, it is convenient to reserve this first section to introduce some notations as well as some standard, yet relevant, concepts regarding these algebras.

Let $V$ and $Z$ be vector spaces over $\mathbb{K} = \mathbb{R}$ or $\mathbb{C}$ and $p \in \mathbb{N}$ with $p \ge 2$. We recall that a $p$-linear map $T: \overbrace{V\times \cdots \times V}^{\text{$p$ times}} \to Z$ is called $p$-alternating if
\begin{eqnarray}
T(\varphi_{1},...,\varphi_{i},...,\varphi_{k},...,\varphi_{p}) = -T(\varphi_{1},...,\varphi_{k},...,\varphi_{i},...,\varphi_{p}) \label{GA1}
\end{eqnarray}
holds for every $\varphi_{1},...,\varphi_{i},...,\varphi_{k},...,\varphi_{p}\in V$. 

\begin{defi}\label{exteriorpower} A $p$-alternating map $\wedge^{p}: V \times \cdots \times V \to Z$ is said to satisfy the universal property for alternating maps if given another vector space $W$ and a $p$-alternating map $T: V \times \cdots \times V \to W$, there exists a unique linear map $\wp : Z \to W$ such that $T = \wp \circ \wedge^{p}$.

Under these conditions, the ordered pair $(Z,\wedge^{p})$ is called a $p$-th exterior power of $V$.
\end{defi}

If $(Z,\wedge^{p})$ is a $p$-th exterior power of $V$, it is customary to write $Z = \bigwedge^{p}V$. This already embodies the $\wedge^{p}$ dependence, so we refer to $\bigwedge^{p}V$ itself as the $p$-th exterior power of $V$. 

We remark that the universal property implies that decomposable vectors of the form $\wedge^{p}(\varphi_{1},...,\varphi_{p}) \equiv \varphi_{1}\wedge \cdots \wedge \varphi_{p}$ generate $\bigwedge^{p} V$ \cite{Greub}, that is, $\operatorname{Span}\wedge^{p} = \bigwedge^{p}V$. Notice that, because $\bigwedge^{p}$ is $p$-alternating, such decomposable vectors are identically zero whenever $\varphi_{i} = \varphi_{j}$ for at least one pair of indices $i \neq j$. In particular, $\varphi \wedge \varphi \equiv \varphi^{2} = 0$ holds true for every $\varphi \in V$. 

The $p$-th exterior power of $V$ can be concretely realized as the quotient space of $\bigotimes^{p}V := \overbrace{V\otimes \cdots \otimes V}^{\text{$p$ times}}$ with the space of all decomposable tensors $\varphi_{1}\otimes \cdots \otimes \varphi_{p}$ with at least one pair of coinciding entries. Moreover, $\bigwedge^{p}V$ can be naturally identified with the space of $p$-fold alternating tensors by means of the universal property. 

When $V$ is finite-dimensional with basis $\mathcal{V} = \{\psi_{1},...,\psi_{n}\}$, the set $\{\psi_{i_{1}}\wedge \cdots \wedge \psi_{i_{p}}: \hspace{0.1cm} 1 \le i_{1} < \cdots < i_{p}\le n\}$ forms a basis for $\bigwedge^{p}V$. Consequently, $\bigwedge^{p}V = \{0\}$ whenever $p > n$. 

Using the conventions $\bigwedge^{0} V := \mathbb{K}$ and $\bigwedge^{1}V := V$, for every fixed $p, q \in \mathbb{N}$ there exists a bilinear map $\wedge: \bigwedge^{p}V \times \bigwedge^{q}V \to \bigwedge^{p+q}V$ which associates, to each pair $(\varphi_{1},...,\varphi_{p},\phi_{1},...,\phi_{q})$, the element $\varphi_{1}\wedge \cdots \wedge \varphi_{p}\wedge \phi_{1}\wedge \cdots \wedge \phi_{q}$. Hence, the algebraic direct sum
\begin{eqnarray}
\bigwedge V := \bigoplus_{p=0}^{\infty}\bigwedge^{p}V, \label{GA2}
\end{eqnarray}
which is a vector space of dimension $2^{n}$, becomes an algebra over $\mathbb{K}$ when equipped with the product
\begin{eqnarray}
f\wedge g := \sum_{p=0}^{\infty}\bigg{(}\sum_{i+j = p}f_{i}\wedge g_{j}\bigg{)}, \label{GA3}
\end{eqnarray}
with $f = (f_{0},f_{1},...),g = (g_{0},g_{1},...) \in \bigwedge V$. This algebra is referred to as the \textit{Grassmann or Exterior algebra} over $\mathbb{K}$. Note that, by identifying each $(0,0,...,f_{p},0,...)$ with $f_{p}\in \bigwedge^{p}V$, any element $f \in \bigwedge V$ can be written as a polynomial of the form
\begin{eqnarray}
f = f(\Psi) = \sum_{p=0}^{\infty}\sum_{i_{1}<\cdots < i_{p}}\zeta_{i_{1},...,i_{p}}(\psi_{i_{1}}\wedge \cdots \wedge \psi_{i_{p}}), \label{GA4}
\end{eqnarray}
with coefficients $\zeta_{i_{1},...,i_{p}}\in \mathbb{K}$. The term $p=0$ in the sum is understood as a scalar. In short, $\bigwedge V$ is the algebra generated by $1$ and the \textit{fields} $\psi_{1},...,\psi_{n}$, which satisfy anti-commutation relations
\begin{eqnarray}
\psi_{i}\wedge \psi_{j} + \psi_{j}\wedge \psi_{i} = 0, \label{GA5}
\end{eqnarray}
for all $i,j=1,...,n$. For this reason, $1, \psi_{1},...,\psi_{n}$ are usually called the \textit{generators} of the Grassmann algebra. 

We call $f$ an \textit{even} (respectively \textit{odd}) element of the algebra if $\zeta_{i_{1},...,i_{p}} = 0$ for every odd (respectively even) $p$ in expression (\ref{GA4}). Of course, every element of the Grassmann algebra can be uniquely decomposed as a sum of an even and an odd element, since
\begin{eqnarray}
\bigwedge V \cong \bigoplus_{\text{$p$ even}}\bigwedge^{p}V \oplus \bigoplus_{\text{$p$ odd}}\bigwedge^{p}V. \nonumber
\end{eqnarray}

Sometimes we need to take other fields into account. Suppose that $U$ is another vector space over $\mathbb{K}$ with basis $\mathcal{U} = \{\theta_{1},...,\theta_{m}\}$, and let us consider the direct sum $V\oplus U$, which has basis $\{(\psi_{1},0),...,(\psi_{n},0),(0,\theta_{1}),...,(0,\theta_{m})\}$. The associated Grassmann algebra $\bigwedge (V\oplus U)$ is generated by $1$ and its basis elements. Under the identifications $(\psi_{i},0) \mapsto \psi_{i}$ and $(0,\theta_{j}) \mapsto \theta_{j}$, a $2$-form $(\psi_{i},0)\wedge (0,\theta_{j})$ becomes simply $\psi_{i}\wedge \theta_{j}$. As a consequence, any element $f \in \bigwedge (V\oplus U)$ can be written as a polynomial
\begin{eqnarray}
f = f(\Psi,\Theta) = \sum_{k+l \le n+m}\sum_{i_{1}<\cdots i_{p}}\sum_{j_{1}<\cdots < j_{l}}\zeta_{i_{1},...,i_{p},j_{1},...,j_{l}}(\psi_{i_{1}}\wedge \cdots \wedge \psi_{i_{p}}\wedge \theta_{j_{1}}\wedge \cdots \wedge \theta_{j_{l}}). \label{GA6}
\end{eqnarray}
In (\ref{GA6}), it is understood that $f$ has only $\theta$-variables when $k=0$ and only $\psi$-variables when $l=0$. Therefore, this expression provides a natural way of identifying both $\bigwedge V$ and $\bigwedge U$ as vector subspaces of $\bigwedge (V\oplus U)$. 

\begin{defi}\label{homogeneouselement} An element $f \in \bigwedge (V\oplus U)$ which depends only on the fields $\psi_{1},...,\psi_{n}$ (resp. $\theta_{1},...,\theta_{n}$) is called $\psi$-homogeneous (resp. $\theta$-homogeneous).
\end{defi}

When $n=m$, we can produce nonhomogeneous elements of $\bigwedge (V\oplus U)$ out of homogeneous ones by a simple translation of variables; if $f = f(\Psi)$ is $\psi$-homogeneous, then
\begin{eqnarray}
f(\Psi+\Theta) := f(\psi_{1}+\theta_{1},...,\psi_{n}+\theta_{n}) \label{GA7}
\end{eqnarray}
is nonhomogeneous. 

Let $f = f_{0}+ f_{1} \in \bigwedge V$, with $f_{0} \in \mathbb{C}$ and $f_{1} \in \bigoplus_{p=1}^{\infty}\bigwedge^{p}V$. Let $U \subseteq \mathbb{C}$ be an open set and $F: U \to \mathbb{C}$ be such that $f_{0} \in U$ and all derivatives $F^{(k)}(f_{0})$ of $F$ exist at $f_{0}$, for $k \le n$. Then, we can define $F(f) \equiv F(f(\Psi)) \in \bigwedge V$ by its Taylor series
\begin{eqnarray}
F(f) := F(f_{0}) + \sum_{k=1}^{\infty}\frac{F^{(k)}(f_{0})}{k!}\;{f_1}^{k}. \label{GA8}
\end{eqnarray}
In particular, we use this result to define the exponential $e^{f}$ for all $f \in \bigwedge V$ and the logarithm $\ln f$ for all those $f \in \bigwedge V$ that satisfy $f_0 > 0$. 

\subsection{Derivatives and Integrals on Grassman Algebras}\label{derivativesandintegrals}

To introduce derivatives on Grassmann algebras, we mimic the rules of partial derivatives of multivariable calculus. Suppose $\bigwedge V$ is generated by $1$ and the fields $\psi_{1},...,\psi_{n}$ and let $k \in \{1,...,n\}$ be fixed. We define the (partial) derivative of a monomial $\psi_{i_{1}}\wedge \cdots \wedge \psi_{i_{p}}$, $0 \le p \le n$, with respect to $\psi_{k}$ to be zero if $p=0$ and 
\beq
\frac{\partial}{\partial \psi_{k}}(\psi_{i_{1}}\wedge \cdots \wedge \psi_{i_{p}}) = \delta_{i_{1},k}\psi_{i_{2}}\wedge \psi_{i_{3}}\wedge  \cdots \wedge \psi_{i_{p}} -\delta_{i_{2},k}\psi_{i_{1}}\wedge \psi_{i_{3}}\wedge \cdots \wedge \psi_{i_{p}}+\cdots +(-1)^{p}\delta_{i_{p},k}\psi_{i_{1}}\wedge \psi_{i_{2}}\wedge \cdots \wedge \psi_{i_{p-1}} \label{GA11}
\eeq
otherwise, with $\delta_{n,m}$ denoting a Kronecker delta. We then extend the definition of derivative to $\bigwedge^{p}V$ by linearity.

\begin{defi}\label{derivativeswedge} If $f$ is given by (\ref{GA4}), then
\begin{eqnarray}
\frac{\partial f}{\partial \psi_{k}} := \sum_{p=0}^{\infty}\sum_{i_{1}<\cdots < i_{p}}\zeta_{i_{1},...,i_{p}}\frac{\partial}{\partial \psi_{k}}(\psi_{i_{1}}\wedge \cdots \wedge \psi_{i_{p}}). \label{GA13}
\end{eqnarray}
\end{defi}

Higher order derivatives are defined by iterative applications of partial derivatives
\begin{eqnarray}
\frac{\partial^{k}f}{\partial{\psi_{i_{1}}\cdots \partial \psi_{i_{k}}}} := \frac{\partial}{\partial \psi_{i_{1}}}\cdots \frac{\partial f}{\partial \psi_{i_{k}}}. \label{GA14}
\end{eqnarray}
Curiously, in the context of Grassmann algebras, integrals and derivatives are defined in the same way.

\begin{defi}\label{integrals} The integral of $f \in \bigwedge V$ with respect to the fields $\psi_{j_{1}},...,\psi_{j_{k}}$ is 
\begin{eqnarray}
\int f d\psi_{j_{k}}\wedge \cdots \wedge d\psi_{j_{1}} := \frac{\partial}{\partial \psi_{j_{1}}}\cdots \frac{\partial f}{\partial \psi_{j_{k}}}. \label{GA15}
\end{eqnarray}
Consequently, the integral of a scalar is always zero whilst
\begin{eqnarray}
\int \psi_{m} d\psi_{n} = \delta_{m,n} \quad \mbox{and} \quad \int f d\psi_{1}\wedge \cdots \wedge d\psi_{n} = \zeta_{1,...,n}. \label{GA16}
\end{eqnarray}
\end{defi}

Before we move on, suppose $f = f_{p} \in \bigwedge^{p}V$ is given by
\begin{eqnarray}
f_{p}(x) =  \sum_{i_{1}<\cdots < i_{p}}\zeta_{i_{1},...,i_{p}}(x) (\psi_{i_{1}}\wedge \cdots \wedge \psi_{i_{p}}), \label{GA17}
\end{eqnarray}
where the coefficients $\zeta_{i_{1},...,i_{p}}$ are now differentiable functions of a real-variable $x$. In this case, we say that $f_{p}(x)$ is differentiable with respect to $x$ and define its derivative by 
\begin{eqnarray}
\frac{d f_{p}(x)}{dx} := \sum_{i_{1}<\cdots < i_{p}}\frac{\zeta_{i_{1},...,i_{p}}(x)}{dx} (\psi_{i_{1}}\wedge \cdots \wedge \psi_{i_{p}}). \label{GA18}
\end{eqnarray}
In this case, if $f(x) = \sum_{p=0}^{\infty}f_{p}(x) \in \bigwedge V$, we define its derivative with respect to $x$ by
\begin{eqnarray}
\frac{d f(x)}{dx} := \sum_{p=0}^{\infty}\sum_{i_{1}<\cdots < i_{p}}\frac{\zeta_{i_{1},...,i_{p}}(x)}{dx} (\psi_{i_{1}}\wedge \cdots \wedge \psi_{i_{p}}). \label{GA19}
\end{eqnarray}

Of course, many standard properties such as linearity and the chain rule for derivatives, as well as rules such as integration  by parts for integrals have natural counterparts in the context of Grassmann algebras, when these operations are defined as before. A much more detailed discussion on these topics can be found, e.g. in \cite{berezin2013introduction, MSbook, FKTbook}.

\subsection{Grassmann Gaussian Integrals}\label{gaussianintegrals}

From now on, let us fix $\mathbb{K} = \mathbb{C}$ and denote by $M_{n}(\mathbb{C})$ the space of all $m \times m$ matrices with complex entries. The determinant of a skew-symmetric matrix $A = (A_{ij}) \in M_{m}(\mathbb{C})$ must satisfy the condition $\det A = (-1)^{m}\det A$ and, for this reason, $A$ is singular when $m$ is odd. On the other hand, when $m=2n$ is even, $n \ge 1$, its determinant is the square of a polynomial $\operatorname{Pf}(A)$, called the Pfaffian of $A$, and given by
\begin{eqnarray}
\operatorname{Pf}(A) := \frac{1}{2^{2n}(2n)!}\;\sum_{\sigma \in S_{2n}}\operatorname{sgn}(\sigma)A_{\sigma(1)\sigma(2)}\cdots A_{\sigma(2n-1)\sigma(2n)}, \label{GA21}
\end{eqnarray}
where the sum ranges over all permutations $\sigma$ of the set $\{1,...,2n\}$. 

\begin{defi}
Let $V$ be a finite-dimensional vector space with basis $\{\psi_{1},...,\psi_{2n}\}$, $n \ge 1$, and $A = (A_{ij}) \in M_{2n}(\mathbb{C})$ skew-symmetric. The Grassmann Gaussian integral on $\bigwedge V$ with covariance (or propagator) $A$ is the linear map
\begin{eqnarray}
\bigwedge V \ni  f \mapsto \E_{\mu_{A}}[f] :=\int f(\Psi)d\mu_{A}(\Psi) \in \mathbb{C} \label{GA22}
\end{eqnarray}
determined by its correlations
\begin{eqnarray}
 \int \psi_{i_{1}}\wedge \cdots \wedge \psi_{i_{p}}d\mu_{A}(\Psi) := \operatorname{Pf}[A_{i_{k},i_{l}}]_{1 \le k,l \le p}. \label{GA23}
\end{eqnarray}
Notice that the pfaffian appearing in the right hand side of (\ref{GA23}) is zero when $p$ is an odd number. 
\end{defi}


We are usually interested in the case where $A$ is not only skew-symmetric, but also invertible. When this is the case, one has
\begin{eqnarray}
\E_{\mu_{A}}[f] = \operatorname{Pf}(A)\int f(\Psi)e^{-\frac{1}{2}\langle \Psi, A^{-1}\Psi\rangle} d\psi_{1}\wedge \cdots \wedge d\psi_{2n}. \label{GA24}
\end{eqnarray}
for every $f \in \bigwedge V$. In (\ref{GA24}), we are expressing fields as vectors
\begin{eqnarray}
\Psi = \begin{pmatrix}
\psi_{1} \\
\vdots \\
\psi_{2n}
\end{pmatrix}
\label{GA25}
\end{eqnarray}
and we have introduced an inner product notation
\begin{eqnarray}
\langle \Psi, \Theta \rangle := \Psi^{T}\Theta = \psi_{1}\wedge \theta_{1} + \cdots + \psi_{2n}\wedge \theta_{2n} \label{GA26}
\end{eqnarray}
which is very convenient when dealing with quadratic forms. In particular, a simple completion of squares argument lead to the following formula
\begin{eqnarray}
\E_{\mu_{A}}[e^{\langle \Psi,\Theta\rangle}] = e^{-\frac{1}{2}\langle \Theta, A\Theta\rangle}. \label{GA27}
\end{eqnarray}

\begin{prop} \label{propertiesgaussianmeasure}

\

\textup{(a)} If the weighted Laplacian operator $\Delta_{A}$ is defined by
\begin{equation}
\Delta_{A} := \bigg{\langle}\frac{\partial}{\partial \Psi}, A\frac{\partial}{\partial \Psi}\bigg{\rangle} \equiv \sum_{i,j=1}^{2n}A_{ij}\frac{\partial}{\partial \psi_{i}}\frac{\partial}{\partial \psi_{j}}  \label{GA28}
\end{equation}
then
\begin{equation}
(\mu_{A}*f)(\Psi) := \int f(\Psi+\Theta)d\mu_{A}(\Theta) = e^{\frac{1}{2}\Delta_{A}}f(\Psi)  \label{GA28.1}
\end{equation}
holds. We refer to (\ref{GA28.1}) as the Heat Kernel Formula. 

\textup{(b)}  If $A, B \in M_{2n}(\mathbb{C})$ be both skew-symmetric and invertible, then
\begin{equation}
\E_{\mu_{A+B}}[f] = \E_{\mu_{B}}[\mu_{A}*f]. \label{GA28.2}
\end{equation}
Relation (\ref{GA28.2}) is also known as the Covariance Splitting Formula. 

\end{prop}
\begin{proof} Part \textup{(a)} is proved in \cite{FKTbook,MSbook}. As for part \textup{(b)}, it is not difficult to see that the formula holds for monomials, so it must hold for every element of the Grassmann algebra by the linearity of the integral. 
\end{proof}

In many applications, the covariance of a fermionic theory can be expressed as a block matrix
\begin{eqnarray}
A = \begin{pmatrix}
0 & C \\
-C & 0 
\end{pmatrix}
\label{GA29}
\end{eqnarray}
for some symmetric matrix $C = (C_{ij}) \in M_{n}(\mathbb{C})$, in which case $\operatorname{Pf}(A) = \operatorname{det}C$. This motivates the split of our fields $\{\psi_{1},...,\psi_{2n}\}$ into two sets $\{\psi_{1},...,\psi_{n}\}$ and $\{\bar{\psi}_{1},...,\bar{\psi}_{n}\}$, with $\bar{\psi}_{i} := \psi_{i+n}$. Such transformations yield
\begin{eqnarray}
\frac{1}{2}\langle \Psi, A\Psi \rangle = \langle \bar{\Psi},C\Psi\rangle. \label{GA30}
\end{eqnarray}
When $C$ is invertible, we have
\begin{eqnarray}
A^{-1} = \begin{pmatrix}
0 & -C^{-1} \\
C^{-1} & 0 
\end{pmatrix}
.\label{GA31}
\end{eqnarray}
Using (\ref{GA24}) and (\ref{GA30}), it is natural to define
\begin{eqnarray}
\int f(\bar{\Psi},\Psi)d\mu_{C}(\bar{\Psi},\Psi) := \det C\int f(\bar{\Psi},\Psi)e^{-\langle \bar{\Psi},C^{-1}\Psi\rangle}d\bar{\psi}_{1}\wedge \cdots \wedge d\bar{\psi}_{n}\wedge d\psi_{1}\wedge \cdots \wedge d\psi_{n}, \label{GA32}
\end{eqnarray}
which is referred to as the Grassmann Gaussian integral with covariance $C$. 

\begin{prop}\label{gaussiancorrelationsbarredunbarred} Let $J = \{j_{1},...,j_{p}\}$ and $K = \{k_{1},...,k_{q}\}$ be nonempty subsets of the set $\{1,...,n\}$. Then, with the above notations
\begin{eqnarray}
\int \psi_{j_{1}}\wedge \cdots \wedge \psi_{j_{p}}\wedge \bar{\psi}_{k_{1}}\wedge \cdots \wedge \bar{\psi}_{k_{q}}d\mu_{C}(\bar{\Psi},\Psi) = \begin{cases}
\displaystyle 0 \quad \mbox{if $p \neq q$} \\
\displaystyle \operatorname{det}C_{J\times K} \quad \mbox{if $p=q$} 
\end{cases}
,\label{GA33}
\end{eqnarray}
with $C_{J\times K} :=(C_{ij})_{i\in J, j \in K}$.
\end{prop}

 
\section{Norms and Correlations Bounds}\label{norms}

Given a skew-symmetric invertible matrix $A = (A_{ij}) \in M_{2n}(\mathbb{C})$, we define its norm by
\begin{eqnarray}
\|A\| \equiv \|A\|_{1,\infty} := \sup_{i}\sum_{j=1}^{2n}|A_{ij}|. \label{NOR1}
\end{eqnarray}
Of course, the introduction of this norm allows us to obtain bounds for correlations, as follows
\begin{eqnarray}
|\operatorname{Pf}[A_{i_{k},i_{l}}]_{1\le k,l \le p}| = |\E_{\mu_{A}}[\psi_{i_{1}}\wedge \cdots \wedge \psi_{i_{p}}]|  \le \|A\|^{p}. \label{NOR2}
\end{eqnarray} 
Unfortunately, these estimates are not suitable for typical applications. This merely reflects that our hypotheses about the matrix $A$ are still too weak, a result of the fact that important features of a typical fermionic covariance are not being taken into account. To strengthen our hypotheses on the covariance, we revisit the discussion in the work of Brydges and Wright \cite{BryWri} and consider $A$ to be the block matrix (\ref{GA29}) and $C$ to be the difference $C = C^{+} - C^{-}$ between two positive-definite matrices $C^{\pm} \in M_{n}(\mathbb{C})$ such that $\det C \neq 0$. The following result is a restatement of Lemma 2.1 and Proposition 2.2 from their work.

\begin{prop}\label{estimatescorrelations} Let $J = \{i_{1},...,i_{p}\}$ be a nonempty subset of the set $\{1,...,2n\}$ and set
\begin{eqnarray}
\Lambda_{J} := \{i \in \{1,...,n\}: \hspace{0.1cm} \mbox{$i \in J$ or $i + n \in J$}\}. \label{NOR3}
\end{eqnarray}
If we set
\begin{eqnarray}
\Psi_{J} := \psi_{i_{1}}\wedge \cdots \wedge \psi_{i_{p}}, \label{NOR4}
\end{eqnarray}
then
\begin{eqnarray}
|\E_{\mu_{A}}[\Psi_{J}]| \le (4\max_{\pm, i \in \Lambda_{J}}C^{\pm}_{ii})^{\frac{|J|}{2}} \label{NOR5}
\end{eqnarray}
where $|J|$ denotes the cardinality of $J$. 
\end{prop}

Proposition \ref{estimatescorrelations} provides suitable bounds for correlations which are going to play a key role in the remaining of our work. Please note, however, that the assumption that $C$ can be decomposed as a difference $C = C^{+}-C^{-}$ of positive covariances $C^{\pm}$ is not essential to develop the method presented in this paper. This assumption was introduced in the original paper \cite{BryWri} and it simplifies the analysis because it immediately leads to Gram type bounds for $\operatorname{Pf}A_{J \times J} = \det C_{J \times J}$, as stated in Proposition \ref{estimatescorrelations}. The applications we are interested here do satisfy this hypothesis (see relation (\ref{PHI4.5})), and for this reason we conveniently adopted it as well. A truly essential hypothesis is briefly mentioned in Remark \ref{boundremark}. 

Once we have set a way of estimating the contribution of a covariance matrix, the next step is to introduce norms on the underlying Grassmann algebra. For this matter, suppose we are given an element $f$ of $\bigwedge V$ whose representation in terms of its generators is 
\begin{eqnarray}
f(\Psi) = \sum_{p=0}^{\infty}\sum_{i_{1}< \cdots < i_{p}}\zeta_{i_{1},...,i_{p}}(\psi_{i_{1}}\wedge \cdots \wedge \psi_{i_{p}}). \nonumber
\end{eqnarray}
Motivated by \cite{BryWri}, we write this expression more compactly as
\begin{eqnarray}
f(\Psi) = \sum_{J\subset \{1,...,2n\}}\zeta_{J}\Psi_{J}, \label{NOR6}
\end{eqnarray}
where the sum ranges over every ordered subset $J= \{i_{1}<\cdots < i_{p}\}$ of $\{1,...,2n\}$ and, to each such subset, $\zeta_{J} := \zeta_{i_{1},...,i_{p}}$ and $\Psi_{J}$ is given as in (\ref{NOR4}). Under these notations, we also set
\begin{eqnarray}
\frac{\partial}{\partial \Psi_{J}} := \frac{\partial^{|J|}}{\partial \psi_{i_{p}}\cdots \partial \psi_{i_{1}}}, \label{NOR7}
\end{eqnarray}
in such a way that
\begin{eqnarray}
\frac{\partial f(\Psi)}{\partial \Psi_{J}} \equiv \zeta_{J}. \label{NOR8}
\end{eqnarray}

\begin{defi}\label{normsgrassmann} Let $z \in \mathbb{R}$ be fixed. The norm of $f \in \bigwedge V$ is defined to be
\begin{eqnarray}
\|f\|_{z} := \sum_{m=1}^{\infty}f_{m}z^{2m}, \label{NOR9}
\end{eqnarray}
where, for each $m$, $f_{m}$ is given by
\begin{eqnarray}
f_{m} := \sup_{i\in \{1,...,2n\}}\frac{1}{2m}\sum_{J \ni i, |J| = 2m}\bigg{|}\frac{\partial f(\Psi)}{\partial \Psi_{J}}\bigg{|}_{\Psi=0}\bigg{|} \equiv \sup_{i\in \{1,...,2n\}}\frac{1}{2m}\sum_{J \ni i, |J| = 2m}|\zeta_{J}|. \label{NOR10}
\end{eqnarray}
The real number $z$ is called a norm parameter.
\end{defi}

%As stressed in \cite{SW}, let us also remark that, although we are referring to $\|\cdot\|_{z}$ as a norm, it is actually a seminorm on the Grassmann algebra because all constant polynomials and odd elements of this algebra have norm zero. This is only a norm when restricted to the subspace of even elements with zero constant terms. All physical objects addressed in this work belong to this subspace and, for this reason, we talk about a norm instead of a seminorm. 

It is important to note that, although we are strictly following the developments of the theory presented in \cite{BryWri,SW}, Definition \ref{normsgrassmann} is the first point at which our work diverges from theirs, as the parameter $z$ is introduced in the norm only with even powers. Of course, this can be done because all physical objects addressed in this work belong to the subspace of even elements with zero constant field. In truth, this is the reason why we even call $\|\cdot\|_{z}$ a norm, given that it is actually a seminorm in the full Grassmann algebra. There is a deeper reason why it is convenient to define (\ref{NOR9}) with only even powers of $z$, but in order to explain it we need to further develop the theory. We return to this point in Remark \ref{remarkpowerz}, in section \ref{majorantmethod}. 

\section{The Renormalization Group Transformation and Polchinski's Equation}\label{rgtransformation}

Suppose $V$ and $U$ are complex vector spaces with basis given, respectively, by $\mathcal{V} = \{\psi_{1},...,\psi_{2n}\}$ and $\mathcal{U} = \{\theta_{1},...,\theta_{2n}\}$ and let $A = (A_{ij}) \in M_{2n}(\mathbb{C})$ be skew-symmetric and invertible. In addition, let $f$ be an even $\psi$-homogeneous element of $\bigwedge(V\oplus U)$, meaning that it is even when seen as an element of $\bigwedge V$. The basic philosophy of the renormalization group (RG) is to study the \textit{effective action}
\begin{eqnarray}
\tilde{F}(\Psi) := -\log [(\mu_{A}*e^{-f})(\Psi)]\label{RGT1}
\end{eqnarray}
not by performing this integral at once but, instead, by exploring the Covariance Splitting Formula to perform a step-by-step integration. At each step, fluctuations are integrated out as a certain range of energy is considered. As discussed in \cite{BryKen,BryWri}, the problem is then reduced to the study of the RG transformation
\begin{eqnarray}
T_{A}: f \mapsto (T_{A}f)(\Psi) :=  -\log [(\mu_{A}*e^{-f})(\Psi)] \label{RGT2}
\end{eqnarray}
which, as in the bosonic case, satisfies a semi-group property
\begin{eqnarray}
(T_{A_{1}+A_{2}}f)(\Psi) = [T_{A_{1}}(T_{A_{2}}f)](\Psi). \label{RGT3}
\end{eqnarray}
Hence, the RG transformation induces a dynamics for the evolution of $f$, called the \textit{bare action} from now on, as the system changes with different energy scales. Following the notations and ideas in \cite{BryWri}, we would like to implement a continuous scale decomposition of the covariance matrix $A$, in which case the evolution law of the RG transformation is governed by a \textit{flow equation}, that is, a partial differential equation (PDE) referred to as the \textit{Polchinski equation}. Thus, the basic hypothesis of our analysis is that $A$ can be decomposed as an integral
\begin{eqnarray}
A = \int_{t_{0}}^{T}\dot{A}(\tau)d\tau \label{RGT4}
\end{eqnarray}
between a lower scale $t_{0}$ and an upper scale $T > t_{0}$. Here, $\dot{A}(\tau)$ denotes the derivative of $A$ with respect to $\tau$. 

Without loss of generality, fix $t_{0} = 0$, $t \in [0,T]$ and let us write (\ref{RGT4}) as a sum
\begin{eqnarray}
A = A(t) + A_{[t,T]}, \label{RGT5}
\end{eqnarray}
with
\begin{eqnarray}
A(t) := \int_{0}^{t}\dot{A}(\tau)d\tau \quad \mbox{and} \quad A_{[t,T]} := \int_{t}^{T}\dot{A}(\tau)d\tau. \label{RGT6}
\end{eqnarray}
If $A(t)$ is assumed to be a block matrix of the form (\ref{GA29}) satisfying the properties in section \ref{norms}, we may define
\begin{eqnarray}
\sigma_{(s,t)}^{2} := \int_{s}^{t}(4\max_{\pm, i \in \Lambda_{J}}\dot{C}_{ii}^{\pm}(\tau))d\tau, \label{RGT7}
\end{eqnarray}
for $0 \le s \le t$. Note that our definition of $\sigma_{(s,t)}^{2}$ differs from the corresponding object in \cite{BryWri}, as our $\sigma_{(0,t)}^{2}$ is homogeneous of degree one, in the sense that a change $t \mapsto t C(\tau)$ leads to a change $\sigma_{(s,t)}^{2} \mapsto t \sigma_{(s,t)}^{2}$. This is in agreement with the discussion in \cite{SW} why the square of the Gram constant must satisfy the aforementioned homogeneity condition. Additionally, from Proposition \ref{estimatescorrelations} we get
\begin{eqnarray}
|\E_{\mu_{A_{[s,t]}}}[\Psi_{J}]|  \le \bigg{|}\bigg{(}4\max_{\pm, i\in \Lambda_{J}}\int_{s}^{t}\dot{C}^{\pm}_{ii}(\tau)d\tau \bigg{)}\bigg{|}^{\frac{|J|}{2}} \le \bigg{|}\int_{s}^{t}(4\max_{\pm,i\in \Lambda_{J}}\dot{C}^{\pm}_{ii}(\tau)d\tau)\bigg{|}^{\frac{|J|}{2}} = (\sigma_{(s,t)}^{2})^{\frac{|J|}{2}}. \label{RGT8}
\end{eqnarray}

\begin{rem}\label{boundremark} As it shall become clear in Section \ref{majorantmethod}, our method works whenever a uniform Gram bound $\gamma_{\dot{C}(\tau)}$ for $\dot{C}(\tau)$ implies a bound $\delta_{C_{[s,t]}}^{2}$ for $C_{[s,t]}$. This condition is enough to guarantee estimate (\ref{RGT8}), even when the covariance is not $C_{[s,t]}$ cannot be represented as a difference $C^{+} - C^{-}$ of positive matrices. 

\end{rem}


By utilizing the semi-group property of the RG transformation on the decomposition (\ref{RGT5}), we can redirect our attention to the scale-dependent functions
\begin{eqnarray}
(T_{A(t)}f)(\Psi) := \tilde{F}(t,\Psi) \quad \mbox{and} \quad \tilde{\varphi}(t,\Psi) := e^{-\tilde{F}(t,\Psi)} =  (\mu_{A(t)}*e^{-f})(\Psi). \label{RGT9}
\end{eqnarray}

We wish to highlight an important aspect of the RG transformation: it preserves parity. This property holds great significance in the analysis that follows, and thus, we present it as a formal proposition for clarity. Its proof can be found in \cite{FKTbook}, Lemma I.23 and, for this reason, it is omitted here. 

\begin{prop}\label{parity} If $f \in \bigwedge (V\oplus U)$ is even and $\psi$-homogeneous, in the sense of Definition \ref{homogeneouselement}, then $F(t,\Psi)$ is an even element of $\bigwedge V$ for every $t \in [0,T]$.
\end{prop}

The fact that $\tilde{\varphi}$ is given by a convolution allows us to use the Heat Kernel Formula
\begin{eqnarray}
\tilde{\varphi}(t,\Psi) = e^{\frac{1}{2}\Delta_{A(t)}}f(\Psi) \label{RGT10}
\end{eqnarray}
to prove that it satisfies the following \textit{heat equation}
\begin{align}
\label{RG11}
\begin{split}
\begin{cases}
\displaystyle\frac{\partial \tilde{\varphi}}{\partial t} = \frac{1}{2}\Delta_{\dot{A}(t)}\tilde{\varphi} \\
\displaystyle \tilde{\varphi}(0,\Psi) = e^{-f(\Psi)}
\end{cases}
\end{split}
\end{align}
whilst $\tilde{F} = -\log \tilde{\varphi}$ satisfies a nonlinear flow equation
\begin{align}
\label{RGT12}
\begin{split}
\begin{cases}
\displaystyle\frac{\partial \tilde{F}}{\partial t} = \frac{1}{2}\Delta_{\dot{A}(t)}\tilde{F}-\frac{1}{2}\langle \nabla \tilde{F}, \dot{A}(t)\nabla \tilde{F}\rangle  \\
\displaystyle \tilde{F}^{0}(\Psi) := \tilde{F}(0,\Psi) = f(\Psi)
\end{cases}
\end{split}
.
\end{align}
The gradient vector $\nabla \tilde{F}$ in (\ref{RGT12}) is defined in the standard way
\begin{eqnarray}
\nabla \tilde{F} := \begin{pmatrix}
\frac{\partial \tilde{F}}{\partial \psi_{1}} \\
\vdots \\
\frac{\partial \tilde{F}}{\partial \psi_{2n}} 
\end{pmatrix}
.\label{RGT13}
\end{eqnarray}

It is more natural to work with normalized functions $F$ and $\varphi$ instead of the unnormalized ones. This is done by setting
\begin{eqnarray}
\varphi(t,\Psi) := \frac{(\mu_{A(t)}*e^{-f})(\Psi)}{\E_{\mu_{A(t)}}[e^{-f}]} =: e^{-F(t,\Psi)}, \label{RGT14}
\end{eqnarray}
in which case we have $\varphi(t, 0) = 1$ and $F(t,0) = 0$ for every $t$. The following result then is an immediate consequence of our previous analysis.

\begin{theo}\label{polchinskiequation} The normalized effective action $F=F(t,\Psi)$ satisfies the flow equation
\begin{eqnarray}
\frac{\partial F}{\partial t} = \frac{1}{2}\Delta_{\dot{A}(t)}F - \frac{1}{2}\langle \nabla F, \dot{A}(t)\nabla F\rangle -\frac{1}{2}\Delta_{\dot{A}(t)}F\bigg{|}_{\Psi = 0} \label{RGT15}
\end{eqnarray}
with $F^{0}(\Psi) := F(0,\Psi) = f(\Psi)$.
\end{theo}

Equation (\ref{RGT15}) is called \textit{Polchinski's equation} after J. Polchinski \cite{Polch}, who first noticed that this PDE could be the basis for a proof of renormalizability\footnote{To see how this is done for $\phi^{4}$ theories, see e.g. \cite{MSbook,KKS} and references therein.}. In principle, if one solves this equation for $F(t,\Psi)$, one gets all the information about the evolution of the bare action. This is why this approach is sometimes called the \textit{exact renormalization group}. Unfortunately, to find a global solution of this equation is too difficult a task and all one can do is to prove local solvability. The idea behind the Majorant Method, which will be discussed in the next section, is to use a Cauchy-Kowalewski type of argument to control the solution of Polchinski's equation by the solution of an associated Hamilton-Jacobi equation, in close analogy to what is done for bosons in \cite{BryKen,BryWri}.
 
As an even element of $\bigwedge V$, the effective action has norm
\begin{eqnarray}
\|F(t)\|_{z} = \sum_{m=1}^{\infty}F_{m}(t)z^{2m} \label{RGT16}
\end{eqnarray}
for each fixed $t$. In particular, when $t=0$ we recover the bare action
\begin{eqnarray}
\|f\|_{z} \equiv \|F^{0}\|_{z} = \sum_{m=1}^{\infty}F_{m}^{0}z^{2m}, \label{RGT17}
\end{eqnarray}
with $F_{m}^{0} = F_{m}(0)$. 

\begin{prop}\label{coefficientsinequality} For each fixed $t \in [0,T]$ and $k \in \mathbb{N}$, the norm coefficients of the effective action satisfy
\begin{eqnarray}
F_{k}(t) \le \sum_{m=k}^{n}F_{m}^{0}\binom{2m}{2k}\sigma_{(0,t)}^{2(m-k)}+\frac{1}{2}\int_{0}^{t}ds ||\dot{A}(s)|| \sum_{\substack{l,m \\ l+m \ge k+1}} F_{l}(s)F_{m}(s)\Gamma_{l,m}(\sigma_{(s,t)}), \label{RGT18}
\end{eqnarray}
with
\begin{eqnarray}
\Gamma_{l,m}(\xi) := 4lm \sum_{\substack{\textup{$k'$, $k''$ odd} \\k'+k'' = 2k}}\binom{2l-1}{k'}\xi^{2l-1-k'}\binom{2m-1}{k''}\xi^{2m-1-k''}. \label{RGT19}
\end{eqnarray}
\end{prop}

The proof of Proposition \ref{coefficientsinequality} is rather technical and, for this reason, it is done in Appendix \ref{proofofproposition}. Nevertheless, this result plays a central role in the construction of a majorant for the effective action. It is worth mentioning that essentially the same result is demonstrated in the Appendix  of \cite{BryWri}; however, it is convenient to revisit their proof in the present context, since the now even exponents of $z$ and $\sigma_{(0,t)}^{2}$ will play an important role in the development of method. 

\section{The Majorant Method}\label{majorantmethod}

We are finally in a position to state and prove our main result, which establishes conditions for the existence of a majorant for the norm of the effective action. First and foremost, it is necessary to  introduce what we mean by a majorant.

\begin{defi} Let $g = g(t,z)$ be an analytic function on the open disc
\begin{eqnarray}
D_{R}:= \{z \in \mathbb{R}: \hspace{0.1cm} |z| < R\} \nonumber
\end{eqnarray}
for every $t \in [a,b]$, $a< b$ fixed. Let 
\begin{eqnarray}
g(t,z) := \sum_{m=1}^{\infty}g_{m}(t)z^{m} \nonumber
\end{eqnarray}
be a power series expansion of $g$ on $[a,b]\times D_{R}$. Suppose $f = f(t,z)$ is given by a formal power series expansion
\begin{eqnarray}
f(t,z) = \sum_{m=1}^{\infty} f_{m}(t)z^{m} \nonumber
\end{eqnarray}
for every $t \in [a,b]$. Then $g$ is said to be a uniform majorant, or simply a majorant, of $f$ if
\begin{eqnarray}
|f_{n}(t)| \le |g_{n}(t)| \nonumber
\end{eqnarray}
holds for every $n$ and $t \in [a,b]$.
\end{defi}

At this point, it is convenient to introduce the following notation and which will be used extensively from now on. We set
\begin{eqnarray}
\phi(0,z) := \|F^{0}\|_{z} = \sum_{m=1}^{\infty}F_{m}^{0}z^{2m} \label{MM1}
\end{eqnarray}
to represent the norm (\ref{RGT17}) of the bare action of our fermionic field theory. We also state the following auxiliary result.  

\begin{prop}\label{auxiliarymajorant} Let $t \ge 0$ be fixed and suppose $\tilde{\phi} = \tilde{\phi}(s,z;t)$ satisfies the following equation
\beq
\tilde{\phi}(s,z;t) = \frac{1}{2}\phi(0,\sigma_{(0,t)}+z)+\frac{1}{2}\phi(0,\sigma_{(0,t)}-z) + \frac{1}{2}\int_{0}^{s} ds' \|\dot{A}(s')\|\bigg{(}\frac{1}{2}\frac{\partial}{\partial z}\bigg{[}\tilde{\phi}(s',\sigma_{(s',t)}+z)+\tilde{\phi}(s',\sigma_{(s',t)}-z)\bigg{]}\bigg{)}^{2} \label{MM2}
\eeq
for each $s \in [0,t]$. Then $\tilde{\phi}$ is the unique solution of the Hamilton-Jacobi equation
\begin{eqnarray}
\frac{\partial\tilde{\phi}}{\partial s}(s,z;t) -\frac{1}{2}\|\dot{A}(s)\|\bigg{(}\frac{\partial \tilde{\phi}}{\partial z}(s,z;t)\bigg{)}^{2} = 0 \label{MM3}
\end{eqnarray}
with initial condition
\begin{eqnarray}
\tilde{\phi}(0,z;t) = \frac{1}{2}[\phi(0,\sigma_{(0,t)}+z)+\phi(0,\sigma_{(0,t)}-z)]. \label{MM4}
\end{eqnarray}
\end{prop}
\begin{proof} It is an easy application of the Fundamental Theorem of Calculus to check that equations (\ref{MM2}) and (\ref{MM3}) are equivalent.
\end{proof}

Our next theorem, which is the central result of this paper, tells us that, under suitable conditions, the auxiliary function $\tilde{\phi}(s,z;t)$ evaluated at $s=t$ becomes a majorant for the norm of our effective $\|F(t)\|_{z}$, at each fixed $t \in [0,T]$. Before getting to the details, let us observe that the introduction of the function
\begin{eqnarray}
\tau(s) := \int_{0}^{s}\|\dot{A}(s')\|ds' \label{MM5}
\end{eqnarray}
for $s \in [0,t]$ allows us to express equation (\ref{MM3}) in a rescaled reduced way, as follows
\begin{eqnarray}
\frac{\partial \tilde{\phi}}{\partial \tau}(\tau,z) - \frac{1}{2}\bigg{(}\frac{\partial \tilde{\phi}}{\partial z}(\tau,z)\bigg{)}^{2} = 0.  \label{MM6}
\end{eqnarray}



\begin{theo} \label{majorant} \textup{\textbf{(Majorant)}} For small $t > 0$, the function $\phi(t,z) := \tilde{\phi}(t,z;t)$ is analytic near $z = 0$ 
\begin{eqnarray}
\phi(t,z) = \sum_{m=1}^{\infty}\phi_{m}(t)z^{2m}. \label{MM7}
\end{eqnarray}
and it satisfies
\begin{eqnarray}
F_{m}(t) \le \phi_{m}(t), \label{MM8}
\end{eqnarray}
for every $m=1,2,3,...$. In other words, $\phi(t,z)$ is a majorant for (\ref{RGT16}). In particular, these inequalities hold for 
\begin{eqnarray}
|z| < R-\sigma_{(0,t)} \label{MM9}
\end{eqnarray}
provided 
\begin{eqnarray}
\sqrt{\tau(t)}+\sigma_{(0,t)}< R, \label{MM10}
\end{eqnarray}
with
\begin{eqnarray}
R^{-2} := \sup_{m \ge 1}(2m F_{m}(0))^{\frac{1}{m}}. \label{MM11}
\end{eqnarray}
\end{theo}
\begin{proof} The construction of a majorant for the norm of the effective action relies on obtaining upper bounds for its norm coefficients. Hence, the key ingredient is Proposition \ref{coefficientsinequality}, as we already anticipated. 

Observe that
\beq
\begin{split}
\sum_{k=1}^{n}\sum_{m=k}^{n}F_{m}^{0}\binom{2m}{2k}\sigma_{(0,t)}^{2(m-k)}z^{2k} 
&\le 
\sum_{m=1}^{k}F_{m}^{0}\sum_{k=0}^{m}\binom{2m}{2k}\sigma_{(0,t)}^{2(m-k)}z^{2k} 
\\
&= 
\frac{1}{2}\sum_{m=1}^{n}F_{m}^{0}[(\sigma_{(0,t)}^{2}+z)^{2m}+(\sigma_{(0,t)}^{2}-z)^{2m}], 
\end{split}\label{MM12}
\eeq
and similarly
\beq
\begin{split}
&\sum_{k=1}^{n}\sum_{\substack{l,m \\ l+m \ge k+1}} F_{l}(s)F_{m}(s)\Gamma_{l,m}(\sigma_{(s,t)})z^{2k} \\
&\le  \sum_{1 \le l,m \le n}4lm F_{l}(s)F_{m}(s)\sum_{\substack{1 \le k' \le 2l-1 \\ 1 \le k'' \le 2m-1 \\ \textup{$k'$, $k''$ odd}}}\binom{2l-1}{k'}\binom{2m-1}{k''}\sigma_{(s,t)}^{2l-1-k'}z^{k'}\sigma_{(s,t)}^{2m-1-k''}z^{k''}  \\
&= \bigg{(}\sum_{m=1}^{n}2mF_{m}(s)\sum_{\textup{$k$ odd}}\binom{2m-1}{k}\sigma_{(s,t)}^{2m-1-k}z^{k}\bigg{)}^{2} \\
&=\bigg{(}\sum_{m=1}^{n}F_{m}(s) \frac{1}{2}\frac{\partial}{\partial z}\bigg{[}(\sigma_{(s,t)}+z)^{2m}+(\sigma_{(s,t)}-z)^{2m}\bigg{]}\bigg{)}^{2} 
\end{split}
\label{MM13}
\eeq
Thus, if we introduce the auxiliary function
\begin{eqnarray}
\gamma(t,z) = \sum_{m=1}^{n}F_{m}(t)z^{2m}, \nonumber
\end{eqnarray}
relations (\ref{MM12}) and (\ref{MM13}) together with (\ref{RGT18}) imply
\beq
\gamma(t,z) \le \frac{1}{2}\gamma(0,\sigma_{(0,t)}+z)+\frac{1}{2}\gamma(0,\sigma_{(0,t)}-z) + \frac{1}{2}\int_{0}^{t}ds ||\dot{A}(s)|| \bigg{(}\frac{1}{2}\frac{\partial}{\partial z}\bigg{[}\gamma(s, \sigma_{(s,t)}+z)+\gamma(s,\sigma_{(s,t)}-z)\bigg{]}\bigg{)}^{2}\nonumber 
\eeq
If $\phi = \phi(t,z)$ denotes the function $\gamma(t,z)$ for which the above inequality becomes an equality, then it is a majorant for $\|F(t)\|_{z}$ and it satisfies
\beq
\phi(t,z) = \frac{1}{2}\phi(0,\sigma_{(0,t)}+z)+\frac{1}{2}\phi(0,\sigma_{(0,t)}-z) + \frac{1}{2}\int_{0}^{t}ds ||\dot{A}(s)|| \bigg{(}\frac{1}{2}\frac{\partial}{\partial z}\bigg{[}\phi(s, \sigma_{(s,t)}+z)+\phi(s,\sigma_{(s,t)}-z)\bigg{]}\bigg{)}^{2}.\label{MM14}
\eeq
by construction. Moreover, it agrees with $\tilde{\phi}(s,z;t)$ at $s=t$, where the latter is the solution of the Hamilton-Jacobi equation (\ref{MM3}) with initial condition (\ref{MM4}). 



To prove the second part of the theorem, notice that the Hamilton-Jacobi equation (\ref{MM3}) preserves inequalities, in the sense of Lemma \ref{comparisionlemma}. Hence, we obtain sufficient conditions for the existence of the majorant by studying this equation with an initial condition which is an upper bound of $\tilde{\phi}(0,z;t)$, as given by (\ref{MM4}).


By the definition of $R$, we readily see that
\beq
\frac{1}{2}[\phi(0,\sigma_{(0,t)}+z)+\phi(0,\sigma_{(0,t)}-z)] \le \frac{1}{2}\sum_{m=1}^{\infty}R^{-2m}(\sigma_{(0,t)}+z)^{2m}+\frac{1}{2}\sum_{m=1}^{\infty}R^{-2m}(\sigma_{(0,t)}-z)^{2m}. \label{MM15}
\eeq
Now, the sum
\begin{eqnarray}
\sum_{m=1}^{\infty}\frac{R^{-(2m-1)}}{2m-1}[(\sigma_{(0,t)}+z)^{2m-1}+(\sigma_{(0,t)}-z)^{2m-1}] \label{MM16}
\end{eqnarray}
has only even powers of $z$, so it can be incorporated to the right hand side of (\ref{MM15}) as to obtain a new upper bound
\beq
\frac{1}{2}[\phi(0,\sigma(0,t)+z)+\phi(0,\sigma_{(0,t)}-z)] \le -\log(1-R^{-1}\sigma_{(0,t)})-\frac{1}{2}\log\bigg{[}1-\frac{R^{-2}z^{2}}{(1-R^{-1}\sigma_{(0,t)})^{2}}\bigg{]}, \label{MM17}
\eeq
provided $R^{-1}\sigma_{(0,t)} < 1$. Writing 
\begin{eqnarray}
\lambda = \lambda(t) := \frac{R^{-1}}{(1-R^{-1}\sigma_{(0,t)})} \label{MM18}
\end{eqnarray}
for convenience, the function
\begin{eqnarray}
 -\log(1-R^{-1}\sigma_{(0,t)})-\frac{1}{2}\log(1-\lambda^{2}z^{2}) \nonumber
\end{eqnarray}
with $|z| \le \lambda$ is a suitable upper bound for $\tilde{\phi}(0,z;t)$. Thus, we need to study (\ref{MM3}) or, equivalently, (\ref{MM6}) with the latter as our new initial condition. This is done by investigating its associated conservation equation, which is satisfied by the new function $u := \partial \tilde{\phi}/\partial z $, and reads
\begin{eqnarray}
\frac{\partial u}{\partial \tau}(\tau,z) - u(\tau,z) \frac{\partial u}{\partial z}(\tau,z) = 0, \label{MM19}
\end{eqnarray}
 with corresponding initial condition
\begin{eqnarray}
u(0,z) \equiv u_{0}(z) = \frac{\lambda^{2}z}{1-\lambda^{2}z^{2}}. \nonumber
\end{eqnarray}

The solution of this initial value problem is attained by means of the method of the characteristics, which is briefly explained in Appendix \ref{hjconservation}. In the present case, it amounts to invert $z_{0} = z_{0}(\tau,z)$ in the following characteristics equation
\begin{eqnarray}
z = z(\tau) = z_{0}\bigg{(}1-\frac{\lambda^{2}\tau}{1-\lambda^{2}z_{0}^{2}}\bigg{)} \label{MM20}
\end{eqnarray}
and plug it back in $u_{0}(z_{0}) \equiv u(\tau,z)$ to recover the desired solution. The issue is that equation (\ref{MM20}) is cubic in $z_{0}$, and we might only get a solution $z_{0} = z_{0}(\tau,z)$ locally. Of course we are interested in solutions that exists around $t = 0$ and, in this regime, $z \sim z_{0}$ is approximately linear and increasing. Consequently, it is sufficient to require that the derivative of $z$ with respect to $z_{0}$ is strictly positive, which yields the following existence condition
\begin{eqnarray}
\lambda^{4}z_{0}^{4} - \lambda^{2}z_{0}^{2}(\lambda^{2}\tau+2) +1-\lambda^{2}\tau > 0. \label{MM21}
\end{eqnarray}
Obviously, this relation holds true for every $z_{0}$ when $t = 0$. Moreover, for $t > 0$ it is sufficient to require that 
\begin{eqnarray}
z_{0}^{2} < \frac{1-\lambda^{2}\tau}{\lambda^{2}(\lambda^{2}\tau +2)}, \label{MM22}
\end{eqnarray}
which will ensure local invertibility whenever (\ref{MM10}) is satisfied. 
\end{proof}

\begin{rem}\label{remarkpowerz} Theorem \ref{majorant} makes it clear that the main reason for using even exponents for the parameter $z$ in the norm is to ensure that it naturally couples with $\sigma_{(0,t)}^{2}$ in the estimates for majorant. The use of $\sigma_{(0,t)}^{2}$ with a square exponent is unavoidable, given that the latter needs to be a homogeneous of degree one upper bound for the covariance decaying. It is exactly because fermionic field theories deal with even polynomials of the Grassmann algebra and the continuous RG transformation preserves this parity (as earlier stated in Proposition \ref{parity}) that we are able to couple these parameters $z$ and $\sigma_{(0,t)}^{2}$ in a very efficient way. 
\end{rem}

Let us notice that given a real positive parameter $\lambda$, the change $f \mapsto \lambda f$ leads to a change $\phi(0,z) \mapsto \lambda \phi(0,z)$ in the norm function, which in its turn rescales $R$ accordingly. Because the majorant in this case is analytic in $z$ and $\lambda$, a Taylor expansion of the effective action around $\lambda = 0$ will be absolutely convergent in some Banach space of interactions. This justifies the assertion frequently found in the literature that the effective action is analytic in the bare action and in the fields \cite{brydges1987expansions, BryKen, SW}. 

Also important is the fact that when the sequence $(2m F_{m}(0))_{m \in \mathbb{N}}$ is bounded and increasing, its limit exists and coincides with the $R^{-2m}$. As a consequence, $R$ is identified with the radius of convergence of $\phi(0,z)$. Furthermore, we can always adjust $R$ large enough so that inequality (\ref{MM10}) is satisfied, as both $\tau(t)$ and $\sigma_{(0,t)}^{2}$ are typically increasing functions of $t$ in applications. Not only this allows us to get all the way up to $t = T$, but also it favors (\ref{MM9}), in the sense that it enlarges the domain of analyticity of the majorant. Of course, choosing $R$ is equivalent to find a Banach space of interactions in which the majorant exists and the bounds are uniformly in the volume and in the fields. 



\section{Applications - Quartic Perturbations}\label{quarticapplications}

The bare action is usually thought as a perturbation of the free theory, so it quite natural to consider the case where $f$ is a quartic interaction. This condition is implemented in our method by considering the special case
\begin{eqnarray}
\|f\|_{z} = \phi(0,z) = \alpha z^{4} \label{QP1}
\end{eqnarray}
for some $\alpha > 0$. Once again, we investigate the associated conservation equation (\ref{MM6}) using the method of the characteristics, only this time the initial condition reads
\begin{eqnarray}
u(0,z) \equiv u_{0}(z) = 12\alpha \sigma_{(0,t)}^{2}z + 4\alpha z^{3}. \label{QP2} 
\end{eqnarray}
For this quartic perturbation, the characteristic equation becomes
\begin{eqnarray}
z = z_{0}-u_{0}(z_{0})\tau = z_{0} - 12\alpha \sigma_{(0,t)}^{2}z_{0}\tau - 4\alpha z_{0}^{3}\tau, \label{QP3}
\end{eqnarray}
and, as before, we need to use (\ref{QP3}) to invert $z_{0} = z_{0}(\tau,z)$, so that $u$ is recovered from the initial condition by $u_{0}(z_{0}) = u_{0}(z_{0}(\tau,z)) = u(t,z)$. 


Interestingly, the equation we obtain in the present case is also cubic in $z_{0}$, just like in the case studied in Theorem \ref{majorant} where the initial condition was logarithmic. Inverting a cubic equation is not a trivial task and we use the same strategy as before, i.e. we search for sufficient conditions for local invertibility.  

In parallel with the logarithmic case, $z$ behaves as a linear function for sufficiently small $t$, a regime we expect the majorant to exist. In close analogy to what was done in the previous section, we require the derivative of $z$ with respect to $z_{0}$ to be strictly positive within a given domain, in which case (\ref{QP3}) would be locally invertible. This condition is satisfied provided the inequality 
\begin{eqnarray}
z_{0}^{2} < \frac{1-12\alpha \sigma_{(0,t)}^{2}\tau(t)}{12\alpha \tau(t)}. \label{QP4}
\end{eqnarray}
holds true. Recall that we are seeking for a solution $u(t,z)$ within a given time range $t \in [0,T]$ for some fixed upper scale $T \in (0,+\infty]$, and both parameters $\sigma_{(0,t)}$ and $\tau(t)$ are supposed to attain a maximum at $T$. Therefore, if we choose $\alpha$ sufficiently small so to assure
\begin{eqnarray}
1-12\alpha \sigma_{(0,T)}^{2}\tau(T) > 0, \label{QP6}
\end{eqnarray}
condition (\ref{QP4}) is certainly satisfied for some $z_{0}$, and local invertibility is proven to hold. This is enough to ensure that our majorant exists at least when $z$ is sufficiently small. Of course, we do not expect to get much farther in a setting as general as the one we are addressing at this point and, therefore, we now turn to an explicit model in which quartic interactions are taken into account. 


\subsection{$\Psi^{4}_{d}$ Theory}\label{phi4d}

We want to consider the regularized free covariance of a fermionic field theory in dimension $d >2$, which is given explicitly by
\begin{eqnarray}
\hat{C}_{\Lambda_{0}}(\pp) = \frac{1}{|\pp|^{2}+m^{2}}e^{-\frac{|\pp|^{2}+m^{2}}{\Lambda_{0}^{2}}}, \quad \pp \in \frac{2\pi}{L}\mathbb{Z}^{d} \label{PHI4.1}
\end{eqnarray}
with $\Lambda_{0} > m > 0$. Setting $\Lambda_{s}:= \Lambda_{0}e^{-s}$ for $s \in [0,\infty)$, we decompose the above covariance into scales
\begin{eqnarray}
\hat{C}_{\Lambda_{0}}(\pp) = \int_{0}^{\infty}\frac{d}{ds'}\bigg{[}-\frac{1}{|\pp|^{2}+m^{2}}e^{-\frac{|\pp|^{2}+m^{2}}{\Lambda_{s'}^{2}}}\bigg{]}ds'. \label{PHI4.2}
\end{eqnarray}
Note how (\ref{PHI4.2}) compares to the scale decomposition presented in the abstract theory developed in section \ref{rgtransformation}: in the present case, $t_{0} = 0$ as before and $T = \infty$ plays the role of the upper scale. The renormalization group transformation arises from splitting $\hat{C}_{\Lambda_{0}}(\pp)$ into a sum
\beq
\hat{C}_{\Lambda_{0}}(\pp) = \int_{0}^{t}\frac{d}{ds'}\bigg{[}-\frac{1}{|\pp|^{2}+m^{2}}e^{-\frac{|\pp|^{2}+m^{2}}{\Lambda_{s'}^{2}}}\bigg{]}ds' + \int_{t}^{\infty}\frac{d}{ds'}\bigg{[}-\frac{1}{|\pp|^{2}+m^{2}}e^{-\frac{|\pp|^{2}+m^{2}}{\Lambda_{s'}^{2}}}\bigg{]}ds', \label{PHI4.3}
\eeq
for some $t \in (0,+\infty)$.

We simplify our notations by defining
\begin{eqnarray}
\hat{C}_{s}(\pp) := \frac{1}{|\pp|^{2}+m^{2}}e^{-\frac{|\pp|^{2}+m^{2}}{\Lambda_{s}^{2}}} \label{PHI4.4}
\end{eqnarray}
and also
\begin{eqnarray}
\hat{C}_{s,t}(\pp) := \hat{C}_{s}(\pp) - \hat{C}_{t}(\pp) = \int_{s}^{t}\frac{d}{ds'}\bigg{[}-\frac{1}{|\pp|^{2}+m^{2}}e^{-\frac{|\pp|^{2}+m^{2}}{\Lambda_{s'}^{2}}}\bigg{]}ds' . \label{PHI4.5}
\end{eqnarray}
Under these new notations, we have $\hat{C}_{0}(\pp) = \hat{C}_{\Lambda_{0}}(\pp)$. 

At each scale $s$, the covariance in position space is obtained by Fourier transform
\begin{eqnarray}
C_{s}(\xx-\yy) :=  L^{-d}\sum_{\pp \in \frac{2\pi}{L}\mathbb{Z}^{d}}\hat{C}_{s}(\pp)e^{i\langle \pp,\xx-\yy\rangle}, \label{PHI4.6}
\end{eqnarray}
and $C_{s,t}(\xx-\yy)$ is defined by $C_{s,t}(\xx-\yy) := C_{s}(\xx-\yy)-C_{t}(\xx-\yy)$, which mimics (\ref{PHI4.5}) in position space. 

To make the connections with the Majorant Method, we need to consider the covariance $A_{0}$, which is formally\footnote{We used the word \enquote{formally} here because the underlying Grassmann algebra has a (countably) infinite number of generators. However, this does not play any role in the analysis because we can always get a theory with a finite number generators by introducing a different cutoff function. In fact, take any strictly decreasing smooth function $\chi: [0,+\infty) \to [0,1]$ such that $$\chi(x) = \begin{cases} \displaystyle 0 \quad \mbox{if $x \ge 4$} \\ 1 \quad \mbox{if $x \le 1$} \end{cases}$$ and set $$\hat{C}_{\Lambda_{0}}(\pp) := \frac{1}{|\pp|^{2}+m^{2}}\chi\bigg{(}\frac{|\pp|^{2}+m^{2}}{\Lambda_{0}^{2}}\bigg{)},$$ which is now zero for $|\pp| > 2\Lambda_{0}$. In this case, we get a theory with $\pp$ restricted to a finite box $$ \pp \in \frac{2\pi}{L}\mathbb{Z}^{d} \cap \{|\pp| < 2\Lambda_{0}\},$$ and the underlying Grassmann algebra becomes finitely generated. In any case, relation (\ref{PHI4.10}) is used just to help us relate $A$ with $C$ to get the parameters of the $\sigma_{(0,t)}^{2}$ and $\|\dot{A}(t)\|$ in terms of $C$. } given by the block matrix (\ref{GA29})
\begin{eqnarray}
A_{0} = \begin{pmatrix}
0 & C_{0} \\
-C_{0} & 0 
\end{pmatrix}
.\label{PHI4.7}
\end{eqnarray}
The decomposition of $A_{0}$ as a sum $A_{[0,t]} + A_{[t,\infty]}$ used to define $F(t,\Psi)$ in section \ref{rgtransformation} now comes from the decomposition (\ref{PHI4.3}). In other words, $A_{[0,t]}$ is the block matrix associated to $C_{0,t}$. 

Note that $\dot{C}_{ii}^{\pm}$ in the present context is replaced by $\dot{C}_{s}(0)$, which can be estimated by
\begin{eqnarray}
|\dot{C}_{s}(0)| \le \int_{\mathbb{R}^{d}}d\pp |\dot{\hat{C}}_{s}(\pp)| =: \|\dot{\hat{C}}_{s}\|_{1}. \label{PHI4.8}
\end{eqnarray}
Using 
\begin{eqnarray}
\|\dot{\hat{C}}_{s}\|_{1} = \frac{2}{\Lambda_{s}^{2}}e^{-\frac{m^{2}}{\Lambda_{s}^{2}}}\int e^{-\frac{|\pp|^{2}}{\Lambda_{s}^{2}}}d\pp = 2\pi^{\frac{d}{2}}e^{-\frac{m^{2}}{\Lambda_{s}^{2}}}\Lambda_{s}^{d-2}, \label{PHI4.9}
\end{eqnarray}
we obtain an upper bound for the integral
\begin{eqnarray}
\int_{s}^{t}ds' \|\dot{\hat{C}}_{s'}\|_{1} = 2\pi^{\frac{d}{2}}\int_{s}^{t}ds'e^{-\frac{m^{2}}{\Lambda_{s'}^{2}}}\Lambda_{s'}^{d-2} \le 2\pi^{\frac{d}{2}}\int_{s}^{t}ds'\Lambda_{s'}^{d-2} = \rho_{d}(\Lambda_{s}^{d-2}-\Lambda_{t}^{d-2}), \label{PHI4.10}
\end{eqnarray}
where we introduced
\begin{eqnarray}
\rho_{d} := \frac{2\pi^{\frac{d}{2}}}{d-2}. \label{PHI4.11}
\end{eqnarray}
Hence, the first parameter of our theory is found
\begin{eqnarray}
\sigma_{(s,t)}^{2} = \rho_{d}(\Lambda_{s}^{d-2}-\Lambda_{t}^{d-2}). \label{PHI4.12}
\end{eqnarray}

We still need to obtain the norm $\|\dot{C}_{s}\|_{1,\infty}$. Because $\dot{\hat{C}} = \frac{2}{\Lambda_s^2}\; e^{- \frac{p^2 + m^2}{\Lambda_s^2} }$ is a Gaussian in $p$, $\dot C$ is Gaussian, too, hence nonnegative, so
\beq
\int_{\mathbb{R}^{d}}d\xx |\dot{C}_{s}(\xx)| = 
\int_{\mathbb{R}^{d}}d\xx \dot{C}_{s}(\xx) =
\dot{\hat{C}}_{s}(0) = \frac{2}{\Lambda_{s}^2}e^{-\frac{m^{2}}{\Lambda_{s}^{2}}}, \label{PHI4.13}
\eeq
we get
\begin{eqnarray}
\|\dot{C}_{s}\|_{1,\infty} \equiv \|\dot{C}_{s}\| = \frac{2}{\Lambda_{s}^2}e^{-\frac{m^{2}}{\Lambda_{s}^{2}}}. \label{PHI4.14}
\end{eqnarray}
Note that the latter vanishes exponentially fast as $\frac{m^{2}}{\Lambda_{s}^{2}} \to +\infty$. It is also uniformly bounded in $s$ because
\begin{eqnarray}
\frac{2}{\Lambda_{s}^2} \; e^{-\frac{m^{2}}{\Lambda_{s}^{2}}} \le \frac{2}{m^{2}}\sup_{\xi \ge 0}\xi e^{-\xi} = \frac{2}{m^{2}}.\label{PHI4.15}
\end{eqnarray}


\subsection{Scaling in $\Psi_{4}^{4}$ Theory}\label{scalingphi4}

Because the majorant for the effective action is obtained by solving a Hamilton-Jacobi equation,  one wants to have coefficient functions in this Hamilton-Jacobi equation that are of order one (\enquote{dimensionless} in the RG terminology), so they do not diverge in the limit one wants to take. We scale the $F_{m}(t)$ in (\ref{RGT16}) by a power of $\Lambda_t$, i.e.\ write 
\beq\label{tilF}
F_{m}(t)=\Lambda_{t}^{a+2bm} \tilde F_{m}(t) \; .
\eeq 
Here $a$ and $b$ are real numbers that depend only on the dimension $d$. The effect of scaling these norm coefficients is that both parameters $\sigma_{(0,t)}$ and $\|\dot{A}(t)\|$ will get rescaled, too.  The aim is to find $a$ and $b$ such that, after this rescaling, the coefficients in the Hamilton-Jacobi equation will be bounded uniformly in $t$. The values of $a$ and $b$ that we shall take are the ones suggested by the results from perturbative renormalization \cite{MSbook}. 

This scaling also means that the initial condition is taken as $\tilde F_m (0) = \tilde F_m^0$, where the $\tilde F_m^0$ are given and bounded independently of $\Lambda_0$, for $m \ge 3$. For $\tilde F_2^0$, we make the same hypothesis here, but in a complete treatment, this function may contain counterterm contributions, which depend on $\Lambda_0$. In this work, we set $\tilde F_1^0 = 0$. 

Let us show in more detail how we rescale $\sigma_{(0,t)}$ and $\|\dot{A}(t)\|$ by scaling the norm coefficients of $F(t,\Psi)$. As we know, the latter satisfies Polchinski's equation, so its norm coefficients $F_{m}(t)$ satisfy (\ref{RGT18}) in Proposition \ref{coefficientsinequality}. Let $s \le t$ and 
\beq
\tilde\sigma_{(s,t)}  = \Lambda_s^b \; \sigma_{(s,t)} \; .
\eeq
Insert (\ref{tilF}) into (\ref{RGT18}) and divide by the factor $\Lambda_t^{a+bk}$, to get an inequality for $\tilde F_k (t)$. The first term on the right hand side of (\ref{RGT18}) becomes
\begin{eqnarray}\label{linterm}
\sum_{m=k}^{n}\frac{\Lambda_{0}^{a+2bm}}{\Lambda_{t}^{a+2bk}} \; \binom{2m}{2k}\; \tilde F_{m}^{0}\; \sigma_{(0,t)}^{2(m-k)} 
=
\left(\frac{\Lambda_t}{\Lambda_0}\right)^{-a -2bk} \;
\sum_{m=k}^n 
\binom{2m}{2k}\;
\tilde F_m^0\;
\tilde\sigma_{(0,t)} ^{2 (m-k)}
\end{eqnarray}
By similar rearrangements, and using (\ref{PHI4.14}), the second term on the right hand side of (\ref{RGT18}) becomes
\beq\label{quadterm}
\int_0^t ds\; 
\left(
\frac{\Lambda_t}{\Lambda_s}
\right)^{-a-2bk} \;
e^{-\frac{m^2}{\Lambda_s^2}} \;
\Lambda_s^{a + 2b -2} \; 
\sum_{\substack{l,m \\ l+m \ge k+1}} \tilde F_{l}(s)\; \tilde F_{m}(s)\;
\Gamma_{l,m}(\tilde\sigma_{(s,t)}) \; .
\eeq
For all $k$ for which $a+2bk \le 0$, and all $0 \le s \le t$, $(\frac{\Lambda_t}{\Lambda_s})^{-a-2bk} \le 1$ can be used both in (\ref{linterm}) and (\ref{quadterm}), so that for these $k$,
\beq\label{goodk}
\tilde F_k (t) 
\le
\sum_{m=k}^n 
\binom{2m}{2k}\;
\tilde F_m^0\;
\tilde\sigma_{(0,t)} ^{2 (m-k)}
+
\int_0^t ds\; 
e^{-\frac{m^2}{\Lambda_s^2}} \;
\Lambda_s^{a + 2b -2} \; 
\sum_{\substack{l,m \\ l+m \ge k+1}} \tilde F_{l}(s)\; \tilde F_{m}(s)\;
\Gamma_{l,m}(\tilde\sigma_{(s,t)}) \; .
\eeq
We now specialize to the case $d=4$. Perturbatively, the model with a local quartic interaction is then just renormalizable, i.e.\ rendering the Schwinger functions finite in the limit $\Lambda_0 \to \infty$ requires only three types of counterterms (mass, field strength, and coupling), but new counterterms of these types are required in every order in perturbation theory \cite{KKS,MSbook}. The analysis that proves this also suggests which exponents to choose: we set $a = 4$ and $b=-1$. Then $a+2b-2=0$ and $a+2bk = 4-2k \le 0$ for all $k \ge 2$, so that leaving out $k=1$ allows us to use (\ref{goodk}). By (\ref{PHI4.12}) with $d=4$, the effect of scaling is that $\sigma_{(s,t)}$ gets replaced with 
\beq
\tilde\sigma_{(s,t)} = \Lambda_s^{-1}\; \sigma_{(s,t)} \le \sqrt{\rho_4}
\eeq
and $\|\dot A (s) \|$ gets replaced with
\beq
a(s) = e^{-\frac{m^2}{\Lambda_s^2}} \le 1\; .
\eeq
Thus with this choice of $a$ and $b$, rescaling has made the coefficients in the integral inequality (\ref{goodk}) uniformly bounded in $s$ and $t$. 

To prove analyticity, we cannot deal with individual values of $k$ separately but use the generating function that gives the majorant, i.e.\ sum over $k$. Since the case $k=1$ is not covered by the inequality, we can only consistently sum over $k$ if we also change the right hand side of the equation by leaving out the all terms of the sum where $l=1$ or $m=1$. In a graphical expansion, this amounts to leaving out the {\em two-point insertions}. This procedure, which corresponds to a projection $\cP_{\ge 2}$ on the field algebra where quadratic terms are mapped to zero, is standard as a first step in perturbative studies of the Polchinski equation and in constructive renormalization group. Specifically, the above-mentioned truncation means that we modify Polchinski's equation (\ref{RGT15}) by applying the projection $\cP_{\ge 2}$ to its right hand side. Including the quadratic terms requires a much more detailed analysis: the flow of the terms with $m \le 2$ has to be traced in more detail, and growing terms need to be controlled by appropriately changed initial conditions (renormalization by counterterms). We are not addressing this renormalization procedure in this work; it requires, among others, estimates for more general weighted norms and Taylor expansion arguments. 

The function $\tau(t)$ we had before is now replaced with
\beq
\tilde\tau (t)  = \int_0^t a(s) \; ds \; .
\eeq 
Since $a$ is smooth and $a(s) \ge 0$ for all $s$, $\tilde \tau$ is smooth and strictly increasing in $t$. 
Because $0 \le a(s) \le 1$, $\tau (t) \le t$. The exponential decay of $a(s)$ at large $s$ implies that $\tilde\tau$ is a bounded function:
\beq\label{ttaubou}
\tilde\tau(t) \le \frac{1}{2e} + \min\{ t, \ln \frac{\Lambda_0}{m} \}
\eeq
For $t \le t_0 = \ln \frac{\Lambda_0}{m}$, this bound follows from $\tau (t) \le t$.
For $t > t_0$, split the integration interval into $[0,t_0)$ and $[t_0,t]$, and use the same estimate on the first interval to get
\beq
\tilde\tau (t) \le \ln \frac{\Lambda_0}{m} + \int_{t_0}^t a(s) \; ds 
\eeq
Substituting $s = t_0 + \sigma$ in the integrand gives
\beq
\int_{t_0}^t a(s) \; ds
=
\int_0^{t-t_0} e^{-e^{2\sigma}} \; d\sigma
\le
\int_0^{\infty} e^{-1-2\sigma} \; d\sigma
=\frac{1}{2e} \; .
\eeq
so (\ref{ttaubou}) holds. Equation (\ref{ttaubou}) makes precise the idea that the flow ``essentially stops'' at $t_0$, i.e.\ when $\Lambda_t$ goes below $m$. 

In condition (\ref{QP4}), $\tilde\tau(t)$ and $\tilde\sigma_{(0,t)}$ now appear instead of $\tau(t)$ and $\sigma_{(0,t)}$. The coefficient $\alpha$ of the quartic interaction does not receive a scaling factor. Thus
\begin{eqnarray}
z_{0}^{2} < \frac{1-12 |\alpha| \tilde\tau(t)\tilde\sigma^{2}_{(0,t)}}{12|\alpha| \tilde\tau(t)} \label{PHI44.18}
\end{eqnarray}
In summary,  we have proven that for $|\alpha|  < \frac{1}{12 \rho_4 (1+ \ln \frac{\Lambda_0}{m})}$, the right hand side of (\ref{PHI44.18}) is positive, the majorant converges and hence the generating function obtained from the Polchinski equation (\ref{RGT15}) with $\cP_{\ge 2}$-projected right hand side is analytic in the fields. This holds for general (not only quartic) interactions $\tilde F_m^0$, provided the norm of $\tilde F^0$ is small enough. (Recall that this means that interaction terms of degree $2m > 4$ get inverse powers of $\Lambda_0$, i.e.\ $F_m^0 = \Lambda^{4-2m} \tilde F_m^0$.)

This bound does not allow to take the limit $\Lambda_0 \to \infty$, and we now discuss the reasons for this. At face value, the main reason is that $\tilde\tau(t)$ grows linearly in $t$, i.e.\ logarithmically in $\Lambda$, before levelling off at about $\ln \frac{\Lambda_0}{m}$, and one may wonder whether one can avoid this by doing more careful bounds. Indeed, the estimate $(\frac{\Lambda_t}{\Lambda_s})^{2k-4} \le 1$ is wasteful for $2k-4 > 0$, and a more careful bound (first integrating over $s$ and then cancelling the power of $\Lambda_t$) removes this logarithm for all $k \ge 3$. Thus in an even stronger truncation that fixes $F_2(t)$ to its initial value $F_2^0$, analyticity can be shown uniformly in $\Lambda_0$. However, for $k=2$ the above bound is an equality, and one can verify in perturbative calculations that this bound is saturated for $k=2$, so that a more detailed analysis is required to determine the signs in the flow of $F_2$, to decide (depending on the sign of $\alpha$ whether there really is a pole when $\alpha \ln\frac{\Lambda_0}{m}$ becomes of order one (known as the Landau pole), or whether asymptotic freedom holds, in which case the upper bound for $\alpha$ is uniform in $\Lambda_0$. 

\section{Concluding Remarks}\label{conclusion}

The purpose of our work was to revisit \cite{BryWri} and to address the gap in their method of controlling Polchinski's equation. For this reason, we followed the latter closely, oftentimes using the same notation and hypothesis, but also taking into account some important remarks in the comments section of \cite{SW}. In particular, we used the same topology in the Grassmann algebra; the only relevant difference is that, in our work, the latter was defined as a power series of the norm parameter $z^{2}$ instead of $z$. This was of course possible because the renormalization group transformation preserves parity, as stated in Proposition \ref{parity}. Although simple, this modification is essential to develop the method. As mentioned in \cite{SW} and reiterated in Section \ref{rgtransformation} of the present paper, the gap in \cite{BryWri} has its roots in their definition of $\sigma_{(s,t)}$, which should be homogeneous of degree one. This is why we considered $\sigma_{(s,t)}^{2}$ here instead of $\sigma_{(s,t)}$. As it turns out, the even powers of $\sigma_{(s,t)}^{2}$ pair naturally with the even powers of $z^{2}$, and the algebra of the norm coefficients can be dealt with to produce the Hamilton-Jacobi (majorant) equation which controls the flow. 

An immediate consequence of the fact that only even exponents in the norm parameter are to be considered in the analysis is that the terms produced by the flow have a backwards-forwards symmetry in $z$. More concretely, both backward $\phi(s,\sigma_{(s,t)}-z)$ and forward $\phi(s,\sigma_{(s,t)}+z)$ translations in $z$ arise naturally in the evolution. Of course, because the method is all about estimating norms, we could have employed further bounds on $\phi(s,\sigma_{(s,t)}-z)$ so to get only contributions for a forward translation $\phi(s,\sigma_{(s,t)}+z)$ instead. This is done in \cite{BryWri}, and represent a second important difference between their work and ours as the use of the Fundamental Theorem of Calculus also allows us to take both translations into consideration. We would expect that the absence of of this further majorization would improve our bounds. However, interestingly enough, apart from a factor of 2 the existence condition expressed by relation (\ref{MM10}) in our Majorant Theorem coincides with the existence condition stated in \cite{BryWri}, even though the objects involved in both conditions have slightly different definitions. On the other hand, these translations also force a $t$ dependence of the initial condition of the majorant equation and it consequently becomes a terminal rather than an initial value problem. 

These observations explain the difference between our majorant equation and the one initially obtained in \cite{BryWri}, where an extra term $\partial\tilde{\phi}/\partial z$ proportional to their covariance decaying parameter $\sigma_{(0,t)}$ is present. Even with this difference, both equations are compatible in the sense that they are both Hamilton-Jacobi equations and their solutions have similar behavior. It is a curious fact that our equation is exactly the same obtained in \cite{BryKen} for bosons. In the latter, the \enquote{wrong sign} in the Hamilton-Jacobi equation makes the Hamiltonian not convex. As explained in \cite{BryKen, BryWri}, this is considered a disavantage in the procedure because, as the majorant evolves according to the flow, it develops singularities immediately when the initial data has not bounded derivatives, which is a too restrictive condition in typical applications. In the case of bosons, this means that we lost relative signs between tree graphs in the expansion of the effective action (or Mayer series, in this case). For fermions, this wrong sign is also a problem, but probably a less restrictive one, since the nilpotency of the Grassmann algebra naturally introduces some sign cancellations in these expansions. This is why perturbation theory of regularized fermionic field theories are convergent.

When an explicit model is being considered and an initial data for $\phi(0,z)$ is provided, it is natural to investigate the majorant equation indirectly by studying its associated conservation equation (\ref{MM19}). This is done by means of the method of characteristics, as we already noted in Section \ref{quarticapplications}. However, the wrong sign problem poses a challenge in this analysis because it leads to crossings of characteristics. In our analysis of the fermionic $\phi^{4}$ model, this is reflected in the fact that we obtain a cubic (thus not globally invertible) characteristic equation. Even with this limitation, we were able to prove existence of the majorant provided $z$ is small enough. 

\begin{appendices}

\section{Proof of Proposition \ref{coefficientsinequality}}\label{proofofproposition}

This section is dedicated to the demonstration of Proposition \ref{coefficientsinequality}. For the reader's convenience, we restate this proposition below.

\begin{prop} For each fixed $t \in [0,T]$ and $k \in \mathbb{N}$, the norm coefficients of the effective action satisfy
\begin{eqnarray}
F_{k}(t) \le \sum_{m=k}^{n}F_{m}^{0}\binom{2m}{2k}\sigma_{(0,t)}^{2(m-k)}+\frac{1}{2}\int_{0}^{t}ds ||\dot{A}(s)|| \sum_{\substack{l,m \\ l+m \ge k+1}} F_{l}(s)F_{m}(s)\Gamma_{l,m}(\sigma_{(s,t)}), \nonumber
\end{eqnarray}
with
\begin{eqnarray}
\Gamma_{l,m}(\xi) := 4lm \sum_{\substack{\textup{$k'$, $k''$ odd} \\k'+k'' = 2k}}\binom{2l-1}{k'}\xi^{2l-1-k'}\binom{2m-1}{k''}\xi^{2m-1-k''}. \nonumber
\end{eqnarray}
\end{prop}
\begin{proof}
Because the unnormalized effective action $\tilde{F}(t,\Psi)$ satisfies the flow equation (\ref{RGT12}), it can be expressed in terms of an integral equation
\begin{eqnarray}
\tilde{F}(t,\Psi) = (\mu_{A(t)}*\tilde{F}^{0})(\Psi) - \frac{1}{2}\int_{0}^{t}ds (\mu_{A_{[s,t]}}*\langle \nabla \tilde{F}(s,\cdot), \dot{A}(s)\nabla \tilde{F}(s,\cdot)\rangle)(\Psi). \label{PRO1}
\end{eqnarray}
Consequently, we have
\beq
\frac{\partial \tilde{F}(t,\Psi)}{\partial \Psi_{J}}\bigg{|}_{\Psi = 0} = \E_{\mu_{A(t)}}\bigg{[}\frac{\partial \tilde{F}^{0}}{\partial \Psi_{J}}\bigg{]} -\frac{1}{2}\int_{0}^{t}ds \sum_{K\subset J}(-1)^{|K|}\E_{\mu_{A_{[s,t]}}}\bigg{[}\bigg{\langle} \nabla \bigg{(}\frac{\partial \tilde{F}(s,\cdot)}{\partial \Psi_{K}}\bigg{)},\dot{A}(s)\nabla \bigg{(}\frac{\partial \tilde{F}(s,\cdot)}{\partial \Psi_{J\setminus K}}\bigg{)}\bigg{\rangle}\bigg{]}, \label{PRO2}
\eeq
by Leibniz rule. Since $F$ an $\tilde{F}$ only differ by a normalization constant, it follows that $F$ must satisfy the same equation, that is
\beq
\frac{\partial F(t,\Psi)}{\partial \Psi_{J}}\bigg{|}_{\Psi = 0} = \E_{\mu_{A(t)}}\bigg{[}\frac{\partial F^{0}}{\partial \Psi_{J}}\bigg{]} -\frac{1}{2}\int_{0}^{t}ds \sum_{K\subset J}(-1)^{|K|}\E_{\mu_{A_{[s,t]}}}\bigg{[}\bigg{\langle} \nabla \bigg{(}\frac{\partial F(s,\cdot)}{\partial \Psi_{K}}\bigg{)},\dot{A}(s)\nabla \bigg{(}\frac{\partial F(s,\cdot)}{\partial \Psi_{J\setminus K}}\bigg{)}\bigg{\rangle}\bigg{]}. \label{PRO3}
\eeq
Notice that we are only interested in the case $|J| = 2k$ has even cardinality. Getting back to the notations introduced in (\ref{NOR6}) and expressing $F^{0}$ by
\begin{eqnarray}
F^{0}(\Psi) = \sum_{J\subset \{1,...,2n\}}\zeta^{0}_{J}\Psi_{J}, \label{PRO4}
\end{eqnarray}
we can readily estimate the first term in the right hand side of (\ref{PRO3})
\begin{eqnarray}
\bigg{|} \E_{\mu_{A(t)}}\bigg{[}\frac{\partial F^{0}}{\partial \Psi_{J}}\bigg{]}\bigg{|} \le |\sum_{K\supset J}\zeta_{K}^{0}\E_{\mu_{A(t)}}[\Psi_{K\setminus J}]| \le \sum_{m=k}^{n}\sum_{\substack{K\supset J \\ |K| = 2m}}|\zeta_{K}^{0}|\sigma_{(0,t)}^{2(m-k)}, \label{PRO5}
\end{eqnarray}
in view of (\ref{RGT8}). Hence,
\beq
\begin{split}
\sup_{i_{0}}\frac{1}{2k}\sum_{J \ni i_{0}, |J| = 2k}\bigg{|} \E_{\mu_{A(t)}}\bigg{[}\frac{\partial F^{0}}{\partial \Psi_{J}}\bigg{]}\bigg{|} &\le \sum_{m=k}^{n}\sup_{i_{0}}\frac{1}{2m}\sum_{K \ni i_{0}, |K| = 2m}|\zeta_{K}^{0}|\frac{2m}{2k}\sum_{\substack{J \subset K, |J| = 2k \\ i_{0} \in J}}\sigma_{(0,t)}^{2(m-k)} \\
&= \sum_{m=k}^{n}F_{m}^{0}\binom{2m}{2k}\sigma_{(0,t)}^{2(m-k)}. 
\end{split}
\label{PRO6}
\eeq
As for the second term in the right hand side of (\ref{PRO3}), we first set
\begin{eqnarray}
F(t,\Psi) = \sum_{J\subset \{1,...,2n\}}\zeta_{J}\Psi_{J}, \label{PRO7}
\end{eqnarray}
in such a way that
\beq
\begin{split}
&\bigg{|}\E_{\mu_{A_{[s,t]}}}\bigg{[}\bigg{\langle} \nabla \bigg{(}\frac{\partial F(s,\cdot)}{\partial \Psi_{K}}\bigg{)},\dot{A}(s)\nabla \bigg{(}\frac{\partial F(s,\cdot)}{\partial \Psi_{J\setminus K}}\bigg{)}\bigg{\rangle}\bigg{]} \bigg{|} 
\\
&\le \sum_{i,j}|\dot{a}_{ij}|\sum_{L \supset K\cup \{i\}}|\zeta_{L}|\sum_{M\supset (J\setminus K)\cup \{j\}}|\zeta_{M}|\E_{\mu_{A_{[s,t]}}}[ \Psi_{L\setminus (K\cup \{i\})}\wedge \Psi_{M\setminus ((J\setminus K)\cup \{j\})}] \\
& \le\sum_{i,j}|\dot{a}_{ij}|\sum_{L\supset K\cup \{i\}}|\zeta_{L}|\sum_{M\supset (J\setminus K)\cup\{j\}}|\zeta_{M}|\sigma_{(s,t)}^{|L|+|M|-|K|-|J\setminus K|-2}\chi_{L\cap M = \emptyset} 
\end{split}
\eeq
where $\chi_{L\cap M=\emptyset}$ indicates that the sum ranges over disjoint sets $L$ and $M$. Summing over $K\subset J$, with $|J| = 2k$, the former is itself majorized by
\beq
\begin{split}
&\sum_{i,j}|\dot{a}_{ij}|\sup_{i_{0}}\frac{1}{2k}\sum_{\substack{J \ni i_{0}\\|J|=2k}}\sum_{K\subset J}\sum_{L\supset K\cup \{i\}}|\zeta_{L}|\sum_{M\supset (J\setminus K)\cup \{j\}}|\zeta_{M}|\sigma_{(s,t)}^{|L|+|M|-|K|-|J\setminus K|-2}\nonumber \\
&\le \sum_{i,j}|\dot{a}_{ij}|\sup_{i_{0}}\frac{1}{2k}\sum_{\substack{J\ni i_{0}\\|J|=2k}}\sum_{K\subset J}\sum_{\substack{l,m\\ l+m\ge k+1}}\sum_{\substack{L\supset K\cup \{i\}\\|L|=2l}}|\zeta_{L}|\sum_{\substack{M\supset (J\setminus K)\cup\{j\}\\|M|=2m}}|\zeta_{M}|\sigma_{(s,t)}^{2l+2m-|K|-|J\setminus K|-2}. \label{PRO8}
\end{split}
\eeq
In (\ref{PRO8}), $i_{0}\in J$ could be an element of either $K$ or $J\setminus K$. Define
\begin{eqnarray}
K' := \begin{cases}
\displaystyle (J\setminus K)\cup \{j\} \quad \mbox{If $i_{0}\in K$} \nonumber \\
\displaystyle K\cup \{i\} \quad \mbox{If $i_{0}\in J\setminus K$} \nonumber
\end{cases}
\end{eqnarray}
If $i_{0} \in K$, (\ref{PRO8}) is further majorized by
\beq
\sup_{j}\sum_{i=1}^{2n}|\dot{a}_{ij}(s)|4lm \sum_{\substack{l,m \\ l+m \ge k+1}}\sigma_{(s,t)}^{2l+2m-2k-2}\sup_{i_{0}}\frac{1}{2l}\sum_{\substack{L\ni i\\|L|=2l}}|\zeta_{L}|\sup_{j_{0}}\frac{1}{2m}\sum_{\substack{M\ni j_{0}\\|M|=2m}}|\zeta_{M}|\sum_{j\in L}\frac{1}{2k}\sum_{\substack{K\subset L\setminus \{i\}\\K'\subset M\\ K\cup K'=J\ni i_{0}\\ |J|=2k}}1 \label{PRO9}
\eeq
If, on the other hand, $i_{0} \in J\setminus K$ then instead of the above majorant, (\ref{PRO8}) is majorized by
\beq
\sup_{i}\sum_{j=1}^{2n}|\dot{a}_{ij}(s)|4lm \sum_{\substack{l,m \\ l+m \ge k+1}}\sigma_{(s,t)}^{2l+2m-2k-2}\sup_{j}\frac{1}{2l}\sum_{\substack{L\ni j\\|L|=2l}}|\zeta_{L}|\sup_{j_{0}}\frac{1}{2m}\sum_{\substack{M\ni j_{0}\\|M|=2m}}|\zeta_{M}|\sum_{i\in L}\frac{1}{2k}\sum_{\substack{K\subset L\setminus \{j\}\\K'\subset M\\ K\cup K'=J\ni i_{0}\\ |J|=2k}}1. \label{PRO10}
\eeq
In either case, both expressions (\ref{PRO9}) and (\ref{PRO10}) are bounded by 
\begin{eqnarray}
||\dot{A}||\sum_{\substack{l,m \\ l+m\ge k+1}} F_{l}(s)F_{m}(s)\Gamma_{l,m}(\sigma_{(s,t)}). \nonumber
\end{eqnarray}
Indeed, when $i_{0} \in K$ we get
\begin{align}
\sum_{\substack{\textup{$k'$,$k''$ odd}\\ k+k'=2k}}\sum_{j\in L}\frac{1}{2(k'+k'')}\sum_{\substack{K\subset L\setminus\{i\}\\|K|=k'\\ K\ni i_{0}}}\sum_{\substack{K'\subset M\\|K'|=k''+1\\K'\ni j}}1 &= \sum_{\substack{\textup{$k'$,$k''$odd}\\k'+k''=2k}}\frac{k''+1}{k'+k''}\binom{2l-1}{k'}\binom{2m-1}{k''} \nonumber \\
&\le \sum_{\substack{\textup{$k'$,$k''$ odd}\\k'+k''=2k}}\binom{2l-1}{k'}\binom{2m-1}{k''} \nonumber
\end{align}
and a similar argument proves this inequality holds when $i_{0} \in J\setminus K$.  
\end{proof}

\section{Hamilton-Jacobi and Conservation Equations - Some Useful Results}\label{hjconservation}

Suppose we are interested in solving the following Hamilton-Jacobi equation
\begin{eqnarray}
\frac{\partial \omega}{\partial t} -\frac{1}{2}\bigg{(}\frac{\partial \omega}{\partial z}\bigg{)}^{2} = 0 \label{HJ1}
\end{eqnarray}
for $t > 0$ and $z \in \mathbb{R}$, with initial condition $\omega(0,z) = g(z)$ for some function $g$. Let us write (\ref{HJ1}) more compactly as
\begin{eqnarray}
\omega_{t}-\frac{1}{2}\omega_{z}^{2} = 0, \label{HJ2}
\end{eqnarray}
where the subscript denotes the derivative with respect to that variable. 

In the theory of Hamilton-Jacobi equations, equation (\ref{HJ1}) is usually expressed as
\begin{eqnarray}
\omega_{t} +H(\omega_{z}) = 0 \nonumber
\end{eqnarray}
with $H(x) := -x^{2}/2$. This function $H$, which is called the Hamiltonian, is typically a convex function of $x$ instead of a concave one, as in the present case. Although most results in the theory of Hamilton-Jacobi apply to convex functions satisfying the additional condition $f(x)/|x| \to \infty$ as $|x| \to \infty$, the following lemma ensures that, even in the case of concave Hamiltonians, equation (\ref{HJ1}) still preserves monotonicity of its solutions.

\begin{lem} \label{comparisionlemma} For $i=1,2$, suppose $\omega_{i} = \omega_{i}(t,z)$ is a solution of the PDE (\ref{HJ1}) with initial condition $\omega_{i}(0,z) = g_{i}(z)$. Then $g_{1} \le g_{2}$ implies $\omega_{1} \le \omega_{2}$. 
\end{lem}
\begin{proof}
Let us define $v_{i} = v_{i}(t,z) := -\omega_{i}(t,z)$, in such a way that it satisfies
\begin{eqnarray}
v_{t} + \frac{1}{2}(v_{z})^{2} = 0. \label{HJE3}
\end{eqnarray}
The solution of (\ref{HJE3}) is given by the Hopf-Lax formula \cite{evans1998partial} 
\begin{eqnarray}
v_{i}(t,z) = \min_{y \in \mathbb{R}}\bigg{\{} \frac{(z-y)^{2}}{t} - g_{i}(y) \bigg{\}}. \nonumber
\end{eqnarray}
Hence, we have
\begin{eqnarray}
\omega_{i}(t,z) = -\min_{y \in \mathbb{R}}\bigg{\{} \frac{(z-y)^{2}}{t} - g_{i}(y) \bigg{\}} = \max_{y \in \mathbb{R}}\bigg{\{} -\frac{(z-y)^{2}}{t} + g_{i}(y) \bigg{\}}, \nonumber
\end{eqnarray}
from where it follows that $\omega_{1} \le \omega_{2}$ when $g_{1} \le g_{2}$. 
\end{proof}

Now, if $\omega = \omega(z,t)$ is at least twice-differentiable and satisfies (\ref{HJ1}), taking derivatives with respect to $z$ yields
\begin{eqnarray}
\omega_{tz} - \omega_{z}\omega_{zz} = 0. \label{HJ4}
\end{eqnarray}
As a consequence, the function $u:= \omega_{z}$ satisfies
\begin{eqnarray}
u_{t} - uu_{z} = 0, \label{HJ5}
\end{eqnarray}
with initial condition $u_{0}(z) := u(0,z) = \omega_{z}(0)$. This is an example of a \textit{conservation equation}. Naturally, the obtaining of a solution of (\ref{HJ5}) allows us to recover the solution $\omega$ of (\ref{HJ1}). 

We want to solve the following initial value problem
\begin{equation}
\begin{cases}
\begin{aligned}
 \displaystyle u_{t} -u u_{z} &= 0  \\
  \displaystyle u(0,z) = u_{0}(z)
  \end{aligned}
  \end{cases}
 ,\label{HJ6}
\end{equation}
which is done here by means of the so-called method of characteristics. Readers unfamiliar with this method can refer to e.g. \cite{evans1998partial}. The first step is to parametrize $z = z(t)$ and define $U(t) := u(t,z(t))$. The characteristic equation becomes
\begin{eqnarray}
z_{t} = -U(t) \label{HJ7}
\end{eqnarray}
with $z(0) := z_{0}$. Notice that
\begin{eqnarray}
U_{t} = u_{t}+u_{z}z_{t} = u_{t} -u u_{z} = 0 \label{HJ8} 
\end{eqnarray}
by hypothesis, from where it follows that
\begin{eqnarray}
U(t) = \operatorname{const} = U(0) = u_{0}(z_{0}). \label{HJ9}
\end{eqnarray}
Putting (\ref{HJ9}) into (\ref{HJ7}), we get
\begin{eqnarray}
z_{t} = z_{0}-u_{0}(z_{0})t. \label{HJ10}
\end{eqnarray}
To solve for $u$, we have to be able to invert $z_{0} = z_{0}(t,z)$ in (\ref{HJ10}). Then, the solution $u = u(t,z)$ is recovered from $u(t,z) = u_{0}(z_{0}(t,z))$. 





\end{appendices}


\noindent
{\bf Acknowledgement. } The authors would like to thank David Brydges for his insightful comments and suggestions on the content of this paper. WK would like to thank Conselho Nacional de Desenvolvimento Cient\'ifico e Tecnol\' ogico (CNPq) for the financial support during his PhD. This study was financed in part by the Coordena\c c\~ ao de Aperfei\c coamento de Pessoal de N\' ivel Superior (CAPES) - Finance Code 001. This work is supported by Deutsche Forschungsgemeinschaft (DFG, German Research Foundation) under Germany's Excellence Strategy  EXC-2181/1 - 390900948 (the Heidelberg STRUCTURES Cluster of Excellence). 

%\bibliography{References}	

\begin{thebibliography}{99}

\bibitem{berezin2012method}
F. Berezin. \reftitel{The method of second quantization.} Elsevier, 2012.

\bibitem{berezin2013introduction}
F. A. Berezin. \reftitel{Introduction to superanalysis.} Springer Science \& Business Media, 2013.

\bibitem{brydges1987expansions}
D.Brydges.
\reftitel{Expansions in statistical mechanics as part of the theory of partial differential equations.} Percolation Theory and Ergodic Theory of Infinite Particle Systems, pages 31-47 (1987)

\bibitem{BryKen}
D.C.\ Brydges and T.\ Kennedy, 
\reftitel{Mayer expansions and the Hamilton-Jacobi equation.} {\em J.\ Stat.\ Phys.\ \bf 48}, 19-49 (1987).

\bibitem{BryWri}
D.C.\ Brydges and J. Wright. 
\reftitel{Mayer expansions and the Hamilton-Jacobi equation II. Fermions, dimensional reduction formulas.} {\em J.\ Stat.\ Phys.\ \bf 51} 435-456 (1988).

\bibitem{BryWriErr}
D.C.\ Brydges and J. Wright. 
\reftitel{Erratum: Mayer expansions and the Hamilton-Jacobi equation II. Fermions, dimensional reduction formulas.} {\em J.\ Stat.\ Phys.\ \bf 97} 1027 (1999).

\bibitem{DeSiKupi}
E.\ De Simone, A.\ Kupiainen, \reftitel{The KAM theorem and renormalization group.} \href{https://arxiv.org/abs/0707.3438}{arXiv:0707.3438}

\bibitem{DR-GN}
M.\ Disertori and V.\ Rivasseau, 
\reftitel{Continuous constructive fermionic renormalization.}
 {\em Annales Henri Poincare \bf 1}, 1-57 (2000)

\bibitem{DR1}
M.\ Disertori and V.\ Rivasseau, 
\reftitel{Interacting Fermi liquid at finite temperature  I: Convergent Attributions.}
 {\em Commun.\ Math.\ Phys.\
\bf 215}, 251  (2000)

\bibitem{DR2}
M.\ Disertori and V.\ Rivasseau, 
\reftitel{Interacting Fermi liquid in two dimensions at finite temperature II: Renormalization.}
 {\em Commun.\ Math.\ Phys.\
\bf 215}, 291  (2000)

\bibitem{Duch}
P.\ Duch, \reftitel{Construction of Gross-Neveu model using Polchinski flow equation.}
arXiv:2403.18562
%https://arxiv.org/abs/2403.18562

\bibitem{evans1998partial}
L. C. Evans. \reftitel{Partial differential equations, volume 19.}  American Mathematical Society, 2022.

\bibitem{FKTbook}
J. S. Feldman, H. Kn\"orrer, and E. Trubowitz. \reftitel{Fermionic functional integrals and the renormalization group,} volume 16. American
Mathematical Society Providence, RI, 2002.

\bibitem{HoStM}
M.B.\ Fr\" ob, J.\ Holland, S.\ Hollands, \reftitel{All-order bounds for correlation functions of gauge-invariant operators in Yang-Mills theory.} {\em J. Math. Phys. \bf 57}, 122301 (2016)

\bibitem{HoHo}
 J.\ Holland, S.\ Hollands, \reftitel{Operator product expansion algebra.} {\em J.Math.Phys. \bf 54} (2013) 072302

\bibitem{KoHoOPE}
S.\ Hollands, C.\ Kopper, \reftitel{The operator product expansion converges in perturbative field theory.} {\em Commun. Math. Phys. \bf 313}, 257-290 (2012)

\bibitem{KeKoQED}
G.\ Keller, C.\ Kopper, \reftitel{Renormalizability proof for QED based on flow equations.} {\em Commun.\ Math.\ Phys.\ \bf 176}, 193-226 (1996)

\bibitem{KeKoZi}
G.\ Keller, C.\ Kopper, \reftitel{Perturbative renormalization of composite operators via flow equations II, Short distance expansion.} {\em Commun. Math. Phys. \bf 153}, 245-276 (1993)

\bibitem{KKS}
G. Keller, C. Kopper, and  M. Salmhofer. \reftitel{Perturbative renormalization and effective Lagrangians in $\phi^4_4$.} {\em Helvetica Physica Acta \bf 65}, 32-52 (1992)

\bibitem{KoMue}
C.\ Kopper, V.F. M\" uller,
\reftitel{Renormalization of spontaneously broken SU(2) Yang-Mills theory with flow equations.}
 {\em Rev.\ Math.\ Phys.\ \bf 21}, 781-820 (2009) 

\bibitem{PeSa}
W.\ Pedra, M.\ Salmhofer,
\reftitel{Determinant bounds and the Matsubara UV problem of many-fermion systems},
{\em Commun. Math. Phys. \bf 282}, 797--818 (2008)

\bibitem{Polch}
J.\ Polchinski, 
\reftitel{Renormalization and effective Lagrangians.}
{\em Nucl.\ Phys.\ \bf B 231}, 269--295 (1984) 

\bibitem{Sa98}
M.\ Salmhofer, \reftitel{Continuous renormalization for fermions
and Fermi liquid theory}, {\em Commun.\ Math.\ Phys.\
\bf 194},  249--295 (1998)

\bibitem{MSbook}
M. Salmhofer. \reftitel{Renormalization: An introduction.} Springer Science \& Business Media, 2013.

\bibitem{SW}
M.\ Salmhofer, C.\ Wieczerkowski, 
\reftitel{Positivity and convergence in fermionic quantum field theory}, {\em J.\ Stat.\ Phys.\ \bf 99}, 557--586 (2000)

\bibitem{Greub}
W. \ Greub. \reftitel{Multilinear Algebra - 2nd Edition} Springer-Verlag, New York-Heidelberg, 1978

\end{thebibliography}


\end{document}

